%  LaTeX support: latex@mdpi.com 
%  For support, please attach all files needed for compiling as well as the log file, and specify your operating system, LaTeX version, and LaTeX editor.

%=================================================================
\documentclass[foundations,article,submit,pdftex,moreauthors]{Definitions/mdpi} 
%\documentclass[preprints,article,submit,pdftex,moreauthors]{Definitions/mdpi} 
% For posting an early version of this manuscript as a preprint, you may use "preprints" as the journal. Changing "submit" to "accept" before posting will remove line numbers.

% Below journals will use APA reference format:
% admsci, aieduc, behavsci, businesses, econometrics, economies, education, ejihpe, famsci, games, humans, ijcs, ijfs, jintelligence, journalmedia, jrfm, jsam, languages, peacestud, psycholint, publications, tourismhosp, youth

% Below journals will use Chicago reference format:
% arts, genealogy, histories, humanities, laws, literature, religions, risks, socsci

%--------------------
% Class Options:
%--------------------
%----------
% journal
%----------
% Choose between the following MDPI journals:
% accountaudit, acoustics, actuators, addictions, adhesives, admsci, adolescents, aerobiology, aerospace, agriculture, agriengineering, agrochemicals, agronomy, ai, aichem, aieduc, aieng, aimater, aimed, aipa, air, aisens, algorithms, allergies, alloys, amh, analog, analytica, analytics, anatomia, anesthres, animals, antibiotics, antibodies, antioxidants, applbiosci, appliedchem, appliedmath, appliedphys, applmech, applmicrobiol, applnano, applsci, aquacj, architecture, arm, arthropoda, arts, asc, asi, astronautics, astronomy, atmosphere, atoms, audiolres, automation, axioms, bacteria, batteries, bdcc, behavsci, beverages, biochem, bioengineering, biologics, biology, biomass, biomechanics, biomed, biomedicines, biomedinformatics, biomimetics, biomolecules, biophysica, bioresourbioprod, biosensors, biosphere, biotech, birds, blockchains, bloods, blsf, brainsci, breath, buildings, businesses, cancers, carbon, cardio, cardiogenetics, cardiovascmed, catalysts, cells, ceramics, challenges, chemengineering, chemistry, chemosensors, chemproc, children, chips, cimb, civileng, cleantechnol, climate, clinbioenerg, clinpract, clockssleep, cmd, cmtr, coasts, coatings, colloids, colorants, commodities, complexities, complications, compounds, computation, computers, condensedmatter, conservation, constrmater, cosmetics, covid, crops, cryo, cryptography, crystals, csmf, ctn, culture, curroncol, cyber, dairy, data, ddc, dentistry, dermato, dermatopathology, designs, devices, dhi, diabetology, diagnostics, dietetics, digital, disabilities, diseases, diversity, dna, drones, dynamics, earth, ebj, ecm, ecologies, econometrics, economies, edm, education, eesp, ejihpe, electricity, electrochem, electronicmat, electronics, encyclopedia, endocrines, energies, eng, engproc, entomology, entropy, environments, environremediat, epidemiologia, epigenomes, esa, est, famsci, fermentation, fibers, fintech, fire, fishes, fluids, foods, forecasting, forensicsci, forests, fossstud, foundations, fractalfract, fuels, future, futureinternet, futurepharmacol, futurephys, futuretransp, galaxies, games, gases, gastroent, gastrointestdisord, gastronomy, gels, genealogy, genes, geographies, geohazards, geomatics, geometry, geosciences, geotechnics, geriatrics, germs, glacies, grasses, green, greenhealth, gucdd, hardware, hazardousmatters, healthcare, hearts, hemato, hematolrep, hep, heritage, higheredu, highthroughput, histories, horticulturae, hospitals, humanities, humans, hydrobiology, hydrogen, hydrology, hydropower, hygiene, idr, iic, ijcs, ijem, ijerph, ijfs, ijgi, ijmd, ijms, ijns, ijom, ijpb, ijt, ijtm, ijtpp, ime, immuno, informatics, information, infrastructures, inorganics, insects, instruments, inventions, iot, j, jaestheticmed, jal, jcdd, jcm, jcp, jcrm, jcs, jcto, jdad, jdb, jdream, jemr, jeta, jfb, jfmk, jgbg, jgg, jimaging, jintelligence, jlpea, jmahp, jmmp, jmms, jmp, jmse, jne, jnt, jof, joi, joitmc, joma, jop, joptm, jor, journalmedia, jox, jpbi, jphytomed, jpm, jrfm, jsam, jsan, jtaer, jvd, jzbg, kidneydial, kinasesphosphatases, knowledge, labmed, laboratories, lae, land, languages, laws, life, lights, limnolrev, lipidology, liquids, literature, livers, logics, logistics, lubricants, lymphatics, machines, macromol, magnetism, magnetochemistry, make, marinedrugs, materials, materproc, mathematics, mca, measurements, medicina, medicines, medsci, membranes, merits, metabolites, metals, meteorology, methane, metrics, metrology, micro, microarrays, microbiolres, microelectronics, micromachines, microorganisms, microplastics, microwave, minerals, mining, mmphys, modelling, molbank, molecules, mps, msf, mti, multimedia, muscles, nanoenergyadv, nanomanufacturing, nanomaterials, ncrna, ndt, network, neuroglia, neuroimaging, neurolint, neurosci, nitrogen, notspecified, nursrep, nutraceuticals, nutrients, obesities, occuphealth, oceans, ohbm, onco, optics, oral, organics, organoids, osteology, oxygen, pandemics, parasites, parasitologia, particles, pathogens, pathophysiology, peacestud, pediatrrep, pets, pharmaceuticals, pharmaceutics, pharmacoepidemiology, pharmacy, philosophies, photochem, photonics, phycology, physchem, physics, physiologia, plants, plasma, platforms, pollutants, polymers, polysaccharides, populations, poultry, powders, precipitation, precisoncol, preprints, proceedings, processes, prosthesis, proteomes, psf, psychiatryint, psychoactives, psycholint, publications, purification, quantumrep, quaternary, qubs, radiation, rdt, reactions, realestate, receptors, recycling, regeneration, religions, remotesensing, reports, reprodmed, resources, rheumato, risks, rjpm, robotics, rsee, ruminants, safety, sci, scipharm, sclerosis, seeds, sensors, separations, sexes, shi, signals, sinusitis, siuj, skins, smartcities, sna, societies, socsci, software, soilsystems, solar, solids, spectroscj, sports, standards, stats, std, stratsediment, stresses, surfaces, surgeries, suschem, sustainability, symmetry, synbio, systems, tae, targets, taxonomy, technologies, telecom, test, textiles, thalassrep, therapeutics, thermo, timespace, tomography, tourismhosp, toxics, toxins, tph, transplantology, transportation, traumacare, traumas, tri, tropicalmed, universe, urbansci, uro, vaccines, vehicles, venereology, vetsci, vibration, virtualworlds, viruses, vision, waste, water, welding, wem, wevj, wild, wind, women, world, youth, zoonoticdis

%---------
% article
%---------
% The default type of manuscript is "article", but can be replaced by: 
% abstract, addendum, article, benchmark, book, bookreview, briefcommunication, briefreport, casereport, changes, clinicopathologicalchallenge, comment, commentary, communication, conceptpaper, conferenceproceedings, correction, conferencereport, creative, datadescriptor, discussion, entry, expressionofconcern, extendedabstract, editorial, essay, erratum, fieldguide, hypothesis, interestingimages, letter, meetingreport, monograph, newbookreceived, obituary, opinion, proceedingpaper, projectreport, reply, retraction, review, perspective, protocol, shortnote, studyprotocol, supfile, systematicreview, technicalnote, viewpoint, guidelines, registeredreport, tutorial,  giantsinurology, urologyaroundtheworld
% supfile = supplementary materials

%----------
% submit
%----------
% The class option "submit" will be changed to "accept" by the Editorial Office when the paper is accepted. This will only make changes to the frontpage (e.g., the logo of the journal will get visible), the headings, and the copyright information. Also, line numbering will be removed. Journal info and pagination for accepted papers will also be assigned by the Editorial Office.

%------------------
% moreauthors
%------------------
% If there is only one author the class option oneauthor should be used. Otherwise use the class option moreauthors.

%---------
% pdftex
%---------
% The option pdftex is for use with pdfLaTeX. Remove "pdftex" for (1) compiling with LaTeX & dvi2pdf (if eps figures are used) or for (2) compiling with XeLaTeX.

%=================================================================
% MDPI internal commands - do not modify
\firstpage{1} 
\makeatletter 
\setcounter{page}{\@firstpage} 
\makeatother
\pubvolume{1}
\issuenum{1}
\articlenumber{0}
\pubyear{2026}
\copyrightyear{2026}
%\externaleditor{Firstname Lastname} % More than 1 editor, please add `` and '' before the last editor name
\datereceived{ } 
\daterevised{ } % Comment out if no revised date
\dateaccepted{ } 
\datepublished{ } 
%\datecorrected{} % For corrected papers: "Corrected: XXX" date in the original paper.
%\dateretracted{} % For retracted papers: "Retracted: XXX" date in the original paper.
%\doinum{} % Used for some special journals, like molbank
%\pdfoutput=1 % Uncommented for upload to arXiv.org
%\CorrStatement{yes}  % For updates
%\longauthorlist{yes} % For many authors that exceed the left citation part
%\IsAssociation{yes} % For association journals

%=================================================================
% Add packages and commands here. The following packages are loaded in our class file: fontenc, inputenc, calc, indentfirst, fancyhdr, graphicx, epstopdf, lastpage, ifthen, float, amsmath, amssymb, lineno, setspace, enumitem, mathpazo, booktabs, titlesec, etoolbox, tabto, xcolor, colortbl, soul, multirow, microtype, tikz, totcount, changepage, attrib, upgreek, array, tabularx, pbox, ragged2e, tocloft, marginnote, marginfix, enotez, amsthm, natbib, hyperref, cleveref, scrextend, url, geometry, newfloat, caption, draftwatermark, seqsplit
% cleveref: load \crefname definitions after \begin{document}

% Commands used by the manuscript.
\newcommand{\pd}[2]{\frac{\partial{#1}}{\partial{#2}}}
\newcommand{\Cc}{\mathbb{C}}
\newcommand{\Ss}{\mathbb{S}}
\newcommand{\pT}{\hat{+}}
\newcommand{\mT}{\hat{-}}
\newcommand{\muT}{\hat{\mu}}
\newcommand{\nuT}{\hat{\nu}}
\newcommand{\muL}{\tilde{\mu}}
\newcommand{\nuL}{\tilde{\nu}}
\newcommand{\uniT}[1]{\mathring{#1}}
\newcommand{\itP}[1]{\hat{#1}}
\newcommand{\lF}[1]{\tilde{#1}}
\newcommand{\Pp}{\mathbb{P}}
\newcommand{\Qq}{\mathbb{Q}}
\def\wh{\hat}
\usepackage{adjustbox}
\usepackage{amsmath}   % For aligned equations
\usepackage{amssymb}   % For math symbols like \mathfrak
\usepackage{booktabs}  % For professional table lines (\toprule, \midrule, etc.)
\usepackage{makecell}  % For forcing line breaks in table headers
\usepackage{placeins}  % For keeping tables with their discussion
\captionsetup{hypcap=false}

%=================================================================
% Please use the following mathematics environments: Theorem, Lemma, Corollary, Proposition, Characterization, Property, Problem, Example, ExamplesandDefinitions, Hypothesis, Remark, Definition, Notation, Assumption
%% For proofs, please use the proof environment (the amsthm package is loaded by the MDPI class).

%=================================================================
% Full title of the paper (Capitalized)
\Title{Interpolating conformal algebra in \texorpdfstring{$(3+1)$}{TEXT} dimensions between the instant form and the light-front form of relativistic dynamics}

% Author Orchid ID: enter ID or remove command
\newcommand{\orcidauthorA}{0000-0002-3276-852X} % Add \orcidA{} behind the author's name
\newcommand{\orcidauthorB}{0000-0002-3024-5186} % Add \orcidB{} behind the author's name

% Authors, for the paper (add full first names)
\Author{Hariprashad Ravikumar $^{1}$\orcidA{} and Chueng-Ryong Ji $^{2}$\orcidB{}}

%\longauthorlist{yes}

% MDPI internal command: Authors, for metadata in PDF
\AuthorNames{Hariprashad Ravikumar and Chueng-Ryong Ji}

%\longauthorlist{yes}

% Affiliations / Addresses (Add [1] after \address if there is only one affiliation.)
\address{%
$^{1}$ \quad Department of Physics, New Mexico State University, Las Cruces, New Mexico 88003-8001, USA; hari1729@nmsu.edu\\
$^{2}$ \quad Department of Physics, North Carolina State University, Raleigh, North Carolina 27695-8202, USA; crji@ncsu.edu}

% Contact information of the corresponding author
\corres{Correspondence: crji@ncsu.edu}

% Current address and/or shared authorship
%\firstnote{Current address: Affiliation.}  
% Current address should not be the same as any items in the Affiliation section.

%\secondnote{These authors contributed equally to this work.}
% The commands \thirdnote{} till \eighthnote{} are available for further notes.

%\simplesumm{} % Simple summary

%\conference{} % An extended version of a conference paper

% Abstract (Do not insert blank lines, i.e. \\) 
\abstract{We extend the interpolation of the Poincar\'e algebra between instant form dynamics (IFD) and light-front dynamics (LFD) to the conformal algebra in $(3+1)$ dimensions. Building upon our recent formulation of the interpolating conformal algebra in $(1+1)$ dimensions, we propose an interpolation method for the conformal group $SO(4,2)$. A central feature of this work is the classification of the $15$ conformal generators into kinematic and dynamic categories as a function of the interpolation angle ($0 \leq \delta \leq \frac{\pi}{4}$). We show that among the five additional generators beyond the Poincar\'e group, the dilatation generator remains strictly kinematic across the entire interpolation region. We also find that in the exact light-front limit ($\delta = \frac{\pi}{4}$), one additional generator from the Special Conformal Transformations (SCT) transforms to become kinematic in $(1+1)$ but becomes dynamical in $(3+1)$. This behavior demonstrates the advantage of LFD in maximizing the number of kinematic generators and minimizing dynamical complexity. To support this algebraic framework, we construct a $6 \times 6$ interpolating projective spacetime matrix representation. We detail the explicit transformations of the interpolating time under all $15$ generators, providing a systematic view of conformal symmetry across the relativistic dynamics.}

% Keywords
\keyword{conformal algebra; light-front dynamics; instant form dynamics; Witt algebra} 

% The fields PACS, MSC, and JEL may be left empty or commented out if not applicable
%\PACS{J0101}
%\MSC{}
%\JEL{}

%%%%%%%%%%%%%%%%%%%%%%%%%%%%%%%%%%%%%%%%%%
% Only for the journal Diversity
%\LSID{\url{http://}}

%%%%%%%%%%%%%%%%%%%%%%%%%%%%%%%%%%%%%%%%%%
% Only for the journal Applied Sciences
%\featuredapplication{Authors are encouraged to provide a concise description of the specific application or a potential application of the work. This section is not mandatory.}
%%%%%%%%%%%%%%%%%%%%%%%%%%%%%%%%%%%%%%%%%%

%%%%%%%%%%%%%%%%%%%%%%%%%%%%%%%%%%%%%%%%%%
% Only for the journal Data
%\dataset{DOI number or link to the deposited data set if the data set is published separately. If the data set shall be published as a supplement to this paper, this field will be filled by the journal editors. In this case, please submit the data set as a supplement.}
%\datasetlicense{License under which the data set is made available (CC0, CC-BY, CC-BY-SA, CC-BY-NC, etc.)}

%%%%%%%%%%%%%%%%%%%%%%%%%%%%%%%%%%%%%%%%%%
% Only for the journal BioTech, Fishes, Neuroimaging and Toxins
%\keycontribution{The breakthroughs or highlights of the manuscript. Authors can write one or two sentences to describe the most important part of the paper.}

%%%%%%%%%%%%%%%%%%%%%%%%%%%%%%%%%%%%%%%%%%
% Only for the journal Encyclopedia
%\encyclopediadef{For entry manuscripts only: please provide a brief overview of the entry title instead of an abstract.}

%%%%%%%%%%%%%%%%%%%%%%%%%%%%%%%%%%%%%%%%%%
% Different journals have different requirements. Please check the specific journal guidelines in the "Instructions for Authors" on the journal's official website.
%\addhighlights{yes}
%\renewcommand{\addhighlights}{%
%
%\noindent The goal is to increase the discoverability and readability of the article via search engines and other scholars. Highlights should not be a copy of the abstract, but a simple text allowing the reader to quickly and simplified find out what the article is about and what can be cited from it. Each of these parts should be devoted up to 2~bullet points.\vspace{3pt}\\
%\textbf{What are the main findings?}
% \begin{itemize}[labelsep=2.5mm,topsep=-3pt]
% \item First bullet.
% \item Second bullet.
% \end{itemize}\vspace{3pt}
%\textbf{What are the implications of the main findings?}
% \begin{itemize}[labelsep=2.5mm,topsep=-3pt]
% \item First bullet.
% \item Second bullet.
% \end{itemize}
%}

%%%%%%%%%%%%%%%%%%%%%%%%%%%%%%%%%%%%%%%%%%
\begin{document}
\section{Introduction}
\label{sec:introduction}

%Conformal symmetry and scope of this work

Studying conformal symmetry is essential for discussing the fundamental properties of spacetime in relativistic quantum field theories~\cite{Wess1960, Kastrup1966, SalamMack1969, Gross1970} and developing characteristics of duality in the anti-de Sitter/conformal field theory (AdS/CFT) correspondence~\cite {Weinberg_2010, Brodsky_2015, Maldacena2016}. In the early 20th century, Cunningham  \cite{Cunningham1910} and Bateman  \cite{Bateman1910} demonstrated that the Maxwell equations are invariant under conformal transformations. Later, Dirac \cite{Dirac1936} proved that the massless version of his equation in relativistic quantum mechanics also exhibits invariance under conformal transformations. In the 1960s, Gross \cite{Gross1970} explored scale invariance as an asymptotic (high-energy) symmetry of scattering amplitudes, demonstrating that the scale invariance of Lagrangian field theories implies conformal invariance. A partial list of literature discussing the physical significance and potential applications of scale and conformal invariance can be found in Refs. \cite{Wess1960, Gürsey1956, Fulton1962, Kastrup1966, SalamMack1969, Wilson1969, Gross1970}. A historical review in Ref. \cite{Kastrup2008} provides further early references. Aspects of the conformal group in two-dimensional Minkowski spacetime, with representations of the algebra and of the group within a field-theoretic Schrodinger representation for bosons and fermions, were reviewed in Ref. \cite{Jackiw1990}.

In our previous work, we investigated the interpolating conformal algebra between instant-form dynamics (IFD) and light-front dynamics (LFD) in $(1+1)$ dimensions \cite{Ji-Hari-2026-PRD}. We constructed a $4\times4$ interpolating projective spacetime matrix representation to map the algebraic structures across all interpolation angles. We explored quantum mechanical operator representations by expressing the conformal generators in terms of the creation and annihilation operators of a simple harmonic oscillator, along with providing a $2\times2$ Pauli matrix representation for the $(1+0)$ and $(0+1)$ conformal groups. We demonstrated that out of the six generators in the $(1+1)$-dimensional conformal algebra, the number of kinematic generators is maximized from two in IFD to four in the LFD limit. This outcome signifies the utility of LFD by proving its capacity to save dynamical efforts in solving $(1+1)$-dimensional relativistic quantum field theories. Building on this foundation, we extend our analysis to $(3+1)$ dimensions in this work.

In general, the conformal transformation $x\longmapsto x'$ in $d$ dimensions can be defined as \cite{Francesco,Blumenhagen,Ji-Hari-2026-PRD},
\begin{align}
\frac{\partial x'^{\alpha}}{\partial x^{\mu}}\frac{\partial x'^{\beta}}{\partial x^{\nu}}g'_{\alpha\beta}=\Lambda(x)g_{\mu\nu},
\end{align}
meaning the metric is preserved up to a scale factor $\Lambda(x)$ under a conformal transformation, where $\mu,\nu\in\{0,\ldots,d-1\}$. The scale $\Lambda(x)=1$ corresponds to the Poincar\'e group consisting of translations, rotations, and Lorentz transformations. Consider an infinitesimal transformation, $x'^{\mu}=x^{\mu}+\epsilon^{\mu}(x)+\mathcal{O}(\epsilon^2)$. The metric then changes by $\delta g_{\mu\nu}=\partial_{\mu}\epsilon_{\nu}(x)+\partial_{\nu}\epsilon_{\mu}(x)$. The conformality condition requires
\begin{align}
\partial_{\mu}\epsilon_{\nu}(x)+\partial_{\nu}\epsilon_{\mu}(x)=F(x)\delta_{\mu\nu},\label{Killing}
\end{align}
where $F(x)=\Lambda(x)-1-\mathcal{O}(\epsilon^2)$. Equation~\eqref{Killing} is the conformal Killing equation. Contraction with $\delta^{\mu\nu}$ yields $F(x)=\frac{2}{d}\partial_{\mu}\epsilon^{\mu}$. There are four classes of solutions for $\epsilon_{\mu}(x)$:
\begin{align}
P_{\mu}&=i\partial_{\mu} &&\text{(translations)},\nonumber\\
M_{\mu\nu}&=i\left(x_{\mu}\partial_{\nu}-x_{\nu}\partial_{\mu}\right) &&\text{(rotations and Lorentz transformations)},\nonumber\\
D&=ix_{\mu}\partial^{\mu} &&\text{(dilatation)},\nonumber\\
\mathfrak{K}_{\mu}&=i\left(2x_{\mu}x_{\nu}\partial^{\nu}-x^2\partial_{\mu}\right) &&\text{(special conformal transformations (SCT))}.
\end{align}
Their corresponding infinitesimal coordinate variations are $\epsilon^{\mu}(x)=a^{\mu}$, $\epsilon^{\mu}(x)=M^{\mu}_{~\nu}x^{\nu}$, $\epsilon^{\mu}(x)=\lambda x^{\mu}$, and $\epsilon^{\mu}(x)=2(b\cdot x)x^{\mu}-x^2b^{\mu}$, respectively. In finite form, the SCT is
\begin{align}
x'^{\mu}=\frac{x^{\mu}-b^{\mu}x^{2}}{1-2b\cdot x+b^{2}x^{2}},
\end{align}
which can be understood as an inversion of $x^\mu$, followed by a translation by $b^\mu$, and then another inversion \cite{Blumenhagen,Ji-Hari-2026-PRD}. These generators form the full conformal algebra, which contains the Poincar\'e group as a subgroup.
\begin{align}\label{conformalalgebra}
    \left[P_{{\mu}},P_{{\nu}}\right]&=0;~\left[\mathfrak{K}_{{\mu}},\mathfrak{K}_{{\nu}}\right]=0;\nonumber\\
 \left[D, P_{{\mu}}\right]&=iP_{{\mu}};~\left[D, \mathfrak{K}_{{\mu}}\right]=-i\mathfrak{K}_{{\mu}};\nonumber\\
 \left[P_{{\rho}},M_{{\mu}{\nu}}\right]&=i\left(g_{{\rho}{\mu}}P_{{\nu}}-g_{{\rho}{\nu}}P_{{\mu}}\right);\nonumber\\
 \left[\mathfrak{K}_{{\rho}},M_{{\mu}{\nu}}\right]&=i\left(g_{{\rho}{\mu}}\mathfrak{K}_{{\nu}}-g_{{\rho}{\nu}}\mathfrak{K}_{{\mu}}\right);\nonumber\\
 \left[M_{{\alpha}{\beta}},M_{{\rho}{\sigma}}\right]&=-i\left(g_{{\beta}{\sigma}}M_{{\alpha}{\rho}}-g_{{\beta}{\rho}}M_{{\alpha}{\sigma}}+g_{{\alpha}{\rho}}M_{{\beta}{\sigma}}-g_{{\alpha}{\sigma}}M_{{\beta}{\rho}}\right);\nonumber\\
 \left[\mathfrak{K}_{{\mu}},P_{{\nu}}\right]&=2i\left(g_{{\mu}{\nu}}D-M_{{\mu}{\nu}}\right);~\left[D, M_{{\mu}{\nu}}\right]=0.
\end{align}
where, $g_{\mu\nu}$ is the metric tensor.

%Poincare interpolation between IFD and LFD

In 1949, for the study of relativistic particle systems, Dirac \cite{Dirac1949} proposed three forms of relativistic dynamics: the instant form ($x^{0}=0$), the front form ($x^{+}=(x^{0}+x^{3})/\sqrt{2}$=0) and the point form ($x^{\mu}x_{\mu}=a^{2}>0, x^{0}>0$).
While the instant form dynamics (IFD) of quantum field theory is produced by quantization at equal time $t=x^{0}$, the light-front dynamics (LFD) is produced by quantization at equal light-front time $\tau \equiv (x^{0}+x^{3})/\sqrt{2}=x^{+}$ ($c=1$ unit is taken here).

One of the main reasons why LFD is advantageous over IFD can be attributed to the energy-momentum dispersion relation. For a particle with mass of $m$ and four-momentum of $k=(k^{0},k^{1},k^{2},k^{3})$, the irrational energy-momentum dispersion relation of the particle at equal-$t$ (IFD) is given by
\begin{align}
  k^{0}=\sqrt{\mathbf{k}^{2}+m^{2}}, \label{eqn:E-P_relation_IF}
\end{align}
where the energy $k^{0}$ is conjugate to $t$ and the three momentum vector $\mathbf{k}$ is given by $\mathbf{k}=(k^{1},k^{2},k^{3})$.
On the contrary, the corresponding rational energy-momentum relation at equal-$\tau$ (LFD) is given by
\begin{align}
  k^{-}=\dfrac{\mathbf{k}_{\perp}^{2}+m^{2}}{k^{+}}, \label{eqn:E-P_relation_LF}
\end{align}
where the light-front energy $k^{-}=(k^{0}-k^{3})/\sqrt{2}$ is conjugate to $\tau$, and the light-front momentum $k^{+}=(k^{0}+k^{3})/\sqrt{2}$ and $\mathbf{k}_{\perp}=(k^{1},k^{2})$ are the orthogonal momenta. This LFD energy-momentum relation is simpler and correlates the signs of $k^{-}$ and $k^{+}$.
When the system evolves in the future direction (positive $\tau$), a positive $k^{-}$ requires a positive $k^{+}$. This feature makes LFD distinct from other forms of relativistic Hamiltonian dynamics and prevents some processes from occurring. For example, spontaneous pair production from the vacuum is forbidden in LFD unless $k^{+}=0$ for both particles due to momentum conservation.

However, this sign correlation does not exist in IFD, and the vacuum structure appears much more complicated than in the LFD case due to the quantum fluctuations. 

\begin{table}[H]
  \caption{\label{tab:Kinematic_and_dynamic_generators_for_different_interoplation_angles}Kinematic and dynamic generators of the Poincaré group for different interpolation angles \cite{Ji2001, Ji2012}}
  \resizebox{\textwidth}{!}{%
    \begin{tabular}{lcc}
	% \hline
	% \hline
	& Kinematic & Dynamic \\
	\hline
	\rule{0pt}{3ex} $\delta=0$ & $\mathcal{K}^{\hat{1}}=-J^{2}, \mathcal{K}^{\hat{2}}=J^{1}, J^{3}, P^{1}, P^{2}, P^{3}$ & $\mathcal{D}^{\hat{1}}=-K^{1}, \mathcal{D}^{\hat{2}}=-K^{2}, -K_3, P^{0}$\\
	$0\leq\delta<\pi/4$ & $\mathcal{K}^{\hat{1}}, \mathcal{K}^{\hat{2}}, J^{3}, P^{1}, P^{2}, P_{\mT}$ & $\mathcal{D}^{\hat{1}}, \mathcal{D}^{\hat{2}}, -K_3, P_{\pT}$\\
	$\delta=\pi/4$ & $\mathcal{K}^{\hat{1}}=-E^{1}, \mathcal{K}^{\hat{2}}=-E^{2}, J^{3}, -K_3, P^{1}, P^{2}, P^{+}$ & $\mathcal{D}^{\hat{1}}=-F^{1}, \mathcal{D}^{\hat{2}}=-F^{2}, P^{-}$\\
	% \hline
	% \hline
      \end{tabular}%
  }
    \end{table}


Even the Poincar\'e algebra is drastically changed in the LFD compared to the IFD.
Among the ten Poincar\'e generators, we have the maximum number (seven) of kinematic (or interaction-independent) operators, which leave the state at $\tau=0$ unchanged. Indeed, the maximum number of kinematic generators allowed in relativistic dynamics is seven. The LFD is the only one that possesses this maximum number of kinematic generators.

LFD maximizes the capacity to
describe hadrons by saving dynamical effort in obtaining quantum chromodynamics (QCD) solutions that reflect the full Poincar\'e symmetries. Determining which generators are kinematic in each form of dynamics is useful, as having more kinematic generators saves dynamical effort when solving relativistic quantum field theories, such as quantum electrodynamics (QED) and QCD~\cite{ji2023relativistic}.

%%%%%%%%%%%%%%
To interpolate the forms of relativistic quantum field theory between IFD and LFD, we take the following interpolating space-time coordinates  \cite{Ji2001, Hornbostel1992, Ji1996, Ji2012, Ji2015EM, Ji2015SP, Ji2018QED, Ji2021QCD} which is defined as a transformation from the ordinary space-time coordinates $x^{\muT}=\mathcal{R}^{\muT}_{\phantom{\mu}{\nu}}x^{\nu}$, i.e.,
\begin{align}\label{eqn:interpolation_angle_definition}
  \begin{pmatrix}
    x^{\hat{+}}\\
    x^{\hat{1}}\\
    x^{\hat{2}}\\
    x^{\hat{-}}
  \end{pmatrix}=
  \begin{pmatrix}
    \cos\delta & 0  & 0  & \sin\delta \\
    0          & 1  & 0  & 0 \\
    0          & 0  & 1  & 0 \\
    \sin\delta & 0  & 0  & -\cos\delta
  \end{pmatrix}
  \begin{pmatrix}
    x^{0}\\
    x^{1}\\
    x^{2}\\
    x^{3}
  \end{pmatrix},
\end{align}
in which the interpolation angle is allowed to run from $0^\circ$ (IFD) through $45^\circ$ (LFD), $0\le \delta \le \frac{\pi}{4}$. The interpolating coordinates $x^{\itP{\pm}}$ in the limit $\delta\rightarrow\pi/4$ become the light-front coordinates $x^{\pm}=(x^{0}\pm x^{3})/\sqrt{2}$ without  ``\textasciicircum''. Note that we interpolate from $-x^3$, so in the limit $\delta\rightarrow\pi/4$, the interpolating coordinates $x^{\itP{-}}$ is $-x^3$ (or $-z$) axis. Since the perpendicular components remain the same ($x^{\itP{j}}=x^{j},x_{\itP{j}}=x_{j}, j=1,2$), we will omit the ``\textasciicircum''  notation unless necessary from now on for the perpendicular indices $j=1,2$ in a four-vector.

 In this interpolating basis, the metric becomes
\begin{align}\label{eqn:g_munu_interpolation}
  g^{\hat{\mu}\hat{\nu}}
  = g_{\hat{\mu}\hat{\nu}}
  =
  \begin{pmatrix}
    \mathbb{C} & 0  & 0  & \mathbb{S} \\
    0 & -1 & 0  & 0 \\
    0 & 0  & -1 & 0 \\
    \mathbb{S} & 0  & 0  & -\mathbb{C}
  \end{pmatrix},
\end{align}
where $\mathbb{S}=\sin2\delta$ and $\mathbb{C}=\cos2\delta$. The lower index variables $x_{\hat{+}}$ and $x_{\hat{-}}$ are related to the upper index variables as $x_{\hat{+}}=g_{\hat{+}\hat{+}}x^{\hat{+}}+g_{\hat{+}\hat{-}}x^{\hat{-}}=\mathbb{C}x^{\hat{+}}+\mathbb{S}x^{\hat{-}}$ and $x_{\hat{-}}=g_{\hat{-}\hat{+}}x^{\hat{+}}+g_{\hat{-}\hat{-}}x^{\hat{-}}=-\mathbb{C}x^{\hat{-}}+\mathbb{S}x^{\hat{+}}$.

The same interpolation
applies to any four-vectors like the four-momentum and the four generators of special conformal transformations. The details of the relationship between the
interpolating variables and the usual space-time variables
can be seen in our previous works  \cite{Ji2001, Hornbostel1992, Ji1996, Ji2012, Ji2015EM, Ji2015SP, Ji2018QED, Ji2021QCD}.

For interpolating the Poincar\'e group, let's start with the Poincar\'e matrix in IFD
\begin{align}\label{eqn:J_mu_nu_IF}
  M_{\mu\nu}=\begin{pmatrix}
     0&-K^1&-K^2&K_3\\
     K^1&0&J^3&-J^2\\
     K^2&-J^3&0&J^1\\
     -K_3&J^2&-J^1&0
     \end{pmatrix}
\end{align}
In interpolation form, it will transform as
\begin{align}\label{eqn:Poincare_Matrix_Interpolation_superscripts}
  M_{\muT\nuT}
  &=\mathcal{R}_{\hat{\mu}}^{{\alpha}}M_{\alpha\beta}\mathcal{R}_{\hat{\nu}}^{{\beta}}=\begin{pmatrix}
    0 & {\mathcal{D}}^{\itP{1}} & {\mathcal{D}}^{\itP{2}} & -K_3\\
    -{\mathcal{D}}^{\itP{1}} & 0 & {J}^{3} & -{\mathcal{K}}^{\itP{1}}\\
    -{\mathcal{D}}^{\itP{2}} & -{J}^{3} & 0 & -{\mathcal{K}}^{\itP{2}}\\
    K_3 & {\mathcal{K}}^{\itP{1}} & {\mathcal{K}}^{\itP{2}} & 0
  \end{pmatrix},
\end{align}
and
\begin{align}
  M^{\muT\nuT}
  =
  g^{\muT\itP{\alpha}}M_{\itP{\alpha}\itP{\beta}}g^{\itP{\beta}\nuT}
  =
  \begin{pmatrix}
    0 & {E}^{\itP{1}} & {E}^{\itP{2}} & K_3\\
    -{E}^{\itP{1}} & 0 & {J}^{3} & -{F}^{\itP{1}}\\
    -{E}^{\itP{2}} & -{J}^{3} & 0 & -{F}^{\itP{2}}\\
    -K_3 & {F}^{\itP{1}} & {F}^{\itP{2}} & 0
  \end{pmatrix}
\end{align}
where
\begin{align}\label{eqn:E_F_D_K_Definition_Interpolation_Angle}
  &E^{\itP{1}}=J^{2}\sin\delta+K^{1}\cos\delta,
  &&\mathcal{K}^{\itP{1}}=-K^{1}\sin\delta-J^{2}\cos\delta, \nonumber\\
  &E^{\itP{2}}=K^{2}\cos\delta-J^{1}\sin\delta,
  &&\mathcal{K}^{\itP{2}}=J^{1}\cos\delta-K^{2}\sin\delta, \nonumber\\
  &F^{\itP{1}}=K^{1}\sin\delta-J^{2}\cos\delta,
  &&\mathcal{D}^{\itP{1}}=-K^{1}\cos\delta+J^{2}\sin\delta, \nonumber\\
  &F^{\itP{2}}=K^{2}\sin\delta+J^{1}\cos\delta,
  &&\mathcal{D}^{\itP{2}}=-J^{1}\sin\delta-K^{2}\cos\delta.
\end{align}
In the limit $\delta=\pi/4$, the interpolating $E^{\itP{j}}$ and $F^{\itP{j}}$ will coincide with the usual $E^{j}$ and $F^{j}$ of LFD.
Note here that the ``\textasciicircum'' notation is reinstated for $j=1, 2$ to emphasize the angle $\delta$ dependence and that the position of the indices on $K, J, E, F, \mathcal{D}, \mathcal{K}$ won't matter, as they are not four-vectors: i.e., $F_{\itP{1}}=F^{\itP{1}}$, etc.

The generalized Poincar\'e algebra for any interpolation angle can be found in \cite{Ji2001}. 
Among the ten Poincar\'e generators, the six generators ($\mathcal{K}^{\itP{1}}, \mathcal{K}^{\itP{2}}, J^{3}, P_{1}, P_{2}, P_{\mT}$) are always kinematic in the sense that the $x^{\pT}=0$ plane remains intact under the transformations generated by them.
As discussed in \cite{Ji2001, Ji2012}, the operator $-K_3=M_{\pT\mT}$ is dynamical in the region where $0\leq\delta<\pi/4$ but becomes kinematic in the light-front limit ($\delta=\pi/4$). Therefore, the instant defined by $x^+=0$ becomes invariant under longitudinal boosts as we move to the light front.
The set of kinematic and dynamic generators depending on the interpolation angle is summarized in Table.~\ref{tab:Kinematic_and_dynamic_generators_for_different_interoplation_angles}.
Since the kinematic transformations do not alter $x^{\pT}$, the individual time-ordered amplitude must be invariant under the kinematic transformations \cite{Ji2018QED}.

In the following sections, we extend this interpolation method to the conformal group, which has the Poincar\'e group as a subgroup. The interpolation method for the conformal group $SO(4,2)$ is proposed in Section~\ref{conformalsimpler}. To characterize the individual properties of the conformal generators, the $6\times6$ interpolating projective spacetime matrix representation is constructed in Section~\ref{sec:projectivespace}, where we systematically classify the kinematic and dynamic generators across the interpolation angles. Finally, a summary and conclusion are presented in Section~\ref{summary-conclusion}.

%\section{Conformal algebra in IFD}
{\color{red} \section{From Poincar\'e Algebra to Conformal Algebra}
\label{sec:conformal-ifd}
In the instant-form limit, we use the un-hatted projective basis
$a,b\in\{-2,-1,0,1,2,3\}$. The fifteen conformal generators are
organized in the following antisymmetric $6\times6$ array \cite{Francesco,Blumenhagen,Ji-Hari-2026-PRD}:
\begin{align}\label{Jab}
  J_{ab}&=
  \begin{pmatrix}
  0&-D&\frac{-\mathfrak{K}_0}{\sqrt{2}}&\frac{-\mathfrak{K}_1}{\sqrt{2}}&\frac{-\mathfrak{K}_2}{\sqrt{2}}&\frac{-\mathfrak{K}_3}{\sqrt{2}}\\
  D&0&\frac{P_0}{\sqrt{2}}&\frac{P_1}{\sqrt{2}}&\frac{P_2}{\sqrt{2}}&\frac{P_3}{\sqrt{2}}\\
    \frac{\mathfrak{K}_0}{\sqrt{2}}&\frac{-P_0}{\sqrt{2}}&0 & -K^{1} & -K^{2} & K_3\\
    \frac{\mathfrak{K}_1}{\sqrt{2}}&\frac{-P_1}{\sqrt{2}}&K^{1} & 0 & J^{3} & -J^{2}\\
    \frac{\mathfrak{K}_2}{\sqrt{2}}&\frac{-P_2}{\sqrt{2}}&K^{2} & -J^{3} & 0 & J^{1}\\
    \frac{\mathfrak{K}_3}{\sqrt{2}}&\frac{-P_3}{\sqrt{2}}&-K_3 & J^{2} & -J^{1} & 0
  \end{pmatrix}_{6\times6}
\end{align}
where $J_{ab}=-J_{ba}$ and $a,b\in\{-2,-1,0,1,2,3\}$. These new generators obey the $SO(4,2)$ commutation relations:
  \begin{align}\label{algebrasimp}
      \left[J_{{a}{b}},J_{{c}{d}}\right]=-i\left(g_{{b}{d}}J_{{a}{c}}-g_{{b}{c}}J_{{a}{d}}+g_{{a}{c}}J_{{b}{d}}-g_{{a}{d}}J_{{b}{c}}\right)
  \end{align}
where, 
  \begin{align}\label{metric}
      g_{ab}=\begin{pmatrix}
  0&-1&0&0&0&0\\
  -1&0&0&0&0&0\\
  0&0&1&0&0&0\\
  0&0&0&-1&0&0\\
  0&0&0&0&-1&0\\
  0&0&0&0&0&-1\\
  \end{pmatrix}_{6\times6}
  \end{align}
The algebra in Eq.~\eqref{algebrasimp}, with the metric in Eq.~\eqref{metric}, is the IFD realization of the conformal algebra. In particular,
$J_{-2,-1}=-D$, $J_{-2,\mu}=-\mathfrak{K}_{\mu}/\sqrt{2}$, and
$J_{-1,\mu}=P_{\mu}/\sqrt{2}$. Its explicit component commutators are collected in Table~\ref{tabelIFD}.

\begingroup
\fontsize{7}{8}\selectfont
\setlength{\tabcolsep}{1pt}
\renewcommand{\arraystretch}{1.25}
\noindent
\begin{minipage}[t]{\textwidth}
\centering
\adjustbox{max width=\textwidth, max totalheight=0.55\textheight}{%
\begin{tabular}{ |c||c|c|c|c|c|c|c|c|c|c|c|c|c|c|c| }
 \hline
 \rule{0pt}{16pt} & $\mathfrak{K}_{{0}}$ & $P_{{0}}$ & $D$ & $\mathfrak{K}_{{3}}$ & $P_{{3}}$ & $K_{{3}}$ & $P_{{1}}$ & $P_{{2}}$ & $\mathfrak{K}_{{1}}$ & $\mathfrak{K}_{{2}}$  & ${K}^{{1}}$ & ${K}^{{2}}$  & ${J}^{{2}}$ & ${J}^{{1}}$ & $J^{{3}}$ \\
 \hline
  \hline
 \rule{0pt}{16pt} $\mathfrak{K}_{{0}}$ &0&$-2iD$&$-i\mathfrak{K}_{{0}}$&0&$2iK_{{3}}$&${i\mathfrak{K}_{{3}}}$&$2i{K}^{{1}}$&$2i{K}^{{2}}$&0&0&$-i\mathfrak{K}_{{1}}$&$-i\mathfrak{K}_{{2}}$&0&0&0\\
 \hline 
 \rule{0pt}{16pt} $P_{{0}}$ &$2iD$&0&$iP_{{0}}$&${2iK_{{3}}}$&0&${iP_{{3}}}$&0&0&$2i{K}^{{1}}$&$2i{K}^{{2}}$&$-iP_{{1}}$&$-iP_{{2}}$&0&0&0\\
 \hline 
 \rule{0pt}{16pt} $D$ &$i\mathfrak{K}_{{0}}$&$-iP_{{0}}$&0&$i\mathfrak{K}_{{3}}$&$-iP_{{3}}$&0&$-iP_{{1}}$&$-iP_{{2}}$&$i\mathfrak{K}_{{1}}$&$i\mathfrak{K}_{{2}}$&0&0&0&0&0\\
 \hline 
 \rule{0pt}{16pt} $\mathfrak{K}_{{3}}$ &0&${-2iK_{{3}}}$&${-i\mathfrak{K}_{{3}}}$&0&$2iD$&${i\mathfrak{K}_{{0}}}$&${-2i{J}^{{2}}}$&${2i{J}^{{1}}}$&0&0&0&0&${-i\mathfrak{K}_{{1}}}$&${i\mathfrak{K}_{{2}}}$&0\\
 \hline 
 \rule{0pt}{16pt} $P_{{3}}$ &$-2iK_{{3}}$&0&$iP_{{3}}$&$-2iD$&0&${iP_{{0}}}$&0&0&${-2i{J}^{{2}}}$&${2i{J}^{{1}}}$&0&0&${-iP_{{1}}}$&${iP_{{2}}}$&0\\
 \hline 
 \rule{0pt}{16pt} $K_{{3}}$ &${-i\mathfrak{K}_{{3}}}$&${-iP_{{3}}}$&0&${-i\mathfrak{K}_{{0}}}$&${-iP_{{0}}}$&0&0&0&0&0&$i{J}^{{2}}$&$-i{J}^{{1}}$&${i{K}^{{1}}}$&$-{i{K}^{{2}}}$&0\\
 \hline 
 \rule{0pt}{16pt} $P_{{1}}$ &$-2i{K}^{{1}}$&0&$iP_{{1}}$&${2i{J}^{{2}}}$&0&0&0&0&$-2iD$&$-2iJ^{{3}}$&$-iP_{{0}}$&0&${iP_{{3}}}$&0&$-iP_{{2}}$\\
 \hline 
 \rule{0pt}{16pt} $P_{{2}}$ &$-2i{K}^{{2}}$&0&$iP_{{2}}$&${-2i{J}^{{1}}}$&0&0&0&0&$2iJ^{{3}}$&$-2iD$&0&$-iP_{{0}}$&0&${-iP_{{3}}}$&$iP_{{1}}$\\
 \hline 
 \rule{0pt}{16pt} $\mathfrak{K}_{{1}}$ &0&$-2i{K}^{{1}}$&$-i\mathfrak{K}_{{1}}$&0&${2i{J}^{{2}}}$&0&$2iD$&$-2iJ^{{3}}$&0&0&$-i\mathfrak{K}_{{0}}$&0&${i\mathfrak{K}_{{3}}}$&0&$-i\mathfrak{K}_{{2}}$\\
 \hline 
 \rule{0pt}{16pt} $\mathfrak{K}_{{2}}$ &0&$-2i{K}^{{2}}$&$-i\mathfrak{K}_{{2}}$&0&${-2i{J}^{{1}}}$&0&$2iJ^{{3}}$&$2iD$&0&0&0&$-i\mathfrak{K}_{{0}}$&0&${-i\mathfrak{K}_{{3}}}$&$i\mathfrak{K}_{{1}}$\\
 \hline 
 \rule{0pt}{16pt} ${K}^{{1}}$ &$i\mathfrak{K}_{{1}}$&$iP_{{1}}$&0&0&0&$-i{J}^{{2}}$&$iP_{{0}}$&0&$i\mathfrak{K}_{{0}}$&0&0&$-iJ^{{3}}$&${iK^{{3}}}$&0&$-i{K}^{{2}}$\\
 \hline 
 \rule{0pt}{16pt} ${K}^{{2}}$ &$i\mathfrak{K}_{{2}}$&$iP_{{2}}$&0&0&0&$i{J}^{{1}}$&0&$iP_{{0}}$&0&$i\mathfrak{K}_{{0}}$&$iJ^{{3}}$&0&0&${-iK^{{3}}}$&$i{K}^{{1}}$\\
 \hline 
 \rule{0pt}{16pt} ${J}^{{2}}$ &0&0&0&${i\mathfrak{K}_{{1}}}$&${iP_{{1}}}$&$-i{K}^{{1}}$&$-iP_{{3}}$&0&${-i\mathfrak{K}_{{3}}}$&0&$-iK^{{3}}$&0&0&$-iJ^{{3}}$&$i{J}^{{1}}$\\
 \hline 
 \rule{0pt}{16pt} ${J}^{{1}}$ &0&0&0&${-i\mathfrak{K}_{{2}}}$&${-iP_{{2}}}$&$i{K}^{{2}}$&0&$iP_{{3}}$&0&${i\mathfrak{K}_{{3}}}$&0&$iK^{{3}}$&$iJ^{{3}}$&0&$-i{J}^{{2}}$\\
 \hline 
 \rule{0pt}{16pt} $J^{{3}}$ &0&0&0&0&0&0&$iP_{{2}}$&$-iP_{{1}}$&$i\mathfrak{K}_{{2}}$&$-i\mathfrak{K}_{{1}}$&$i{K}^{{2}}$&$-i{K}^{{1}}$&$-i{J}^{{1}}$&$i{J}^{{2}}$&0\\
 \hline 
\end{tabular}}
\captionof{table}{Conformal algebra $(3+1)$ in the IFD limit ($\delta \to 0$)}
\label{tabelIFD}
\end{minipage}
\endgroup

%\section{Isomorphism between Dirac matrices and the IFD conformal algebra}

Our previous study of the $(1+1)$-dimensional interpolating conformal
algebra constructed a $4\times4$ interpolating projective representation and
used $2\times2$ Pauli-matrix realizations for the $(1+0)$- and
$(0+1)$-dimensional conformal groups \cite{Ji-Hari-2026-PRD}. Here we give
the corresponding four-dimensional Dirac-matrix realization only for the IFD
conformal algebra. The light-front and general interpolating realizations are
obtained by applying the projective interpolation transformation introduced
below; no independent hatted gamma-matrix convention is required.

The Dirac matrices obey the Clifford algebra
\begin{align}
\{\gamma^\mu,\gamma^\nu\}=2g^{\mu\nu}I_4,\qquad
\gamma_\mu=g_{\mu\nu}\gamma^\nu,\qquad
\gamma_5=i\gamma^0\gamma^1\gamma^2\gamma^3.
\end{align}
The following dictionary realizes all fifteen IFD conformal generators:
\begin{align}
P_\mu&=\frac{i}{2}(1+\gamma_5)\gamma_\mu,
&\mathfrak{K}_\mu&=-\frac{i}{2}(1-\gamma_5)\gamma_\mu,\nonumber\\
D&=\frac{i}{2}\gamma_5,
&K^i&=\frac{i}{2}\gamma_i\gamma_0,\qquad i=1,2,3,\nonumber\\
J^1&=\frac{i}{2}\gamma_2\gamma_3,
&J^2&=\frac{i}{2}\gamma_3\gamma_1,
\qquad J^3=\frac{i}{2}\gamma_1\gamma_2.
\label{eq:ifd-dirac-dictionary}
\end{align}
The longitudinal generator in the projective array has the lowered-index
form $K_3=-K^3$. Substitution of Eq.~\eqref{eq:ifd-dirac-dictionary} into
Eq.~\eqref{Jab} gives the spinor realization
\begingroup\footnotesize
\begin{align}
\rho_\gamma(J_{ab})=\frac{i}{2}
\begin{pmatrix}
0&-\gamma_5&\dfrac{(1-\gamma_5)\gamma_0}{\sqrt{2}}&\dfrac{(1-\gamma_5)\gamma_1}{\sqrt{2}}&\dfrac{(1-\gamma_5)\gamma_2}{\sqrt{2}}&\dfrac{(1-\gamma_5)\gamma_3}{\sqrt{2}}\\
\gamma_5&0&\dfrac{(1+\gamma_5)\gamma_0}{\sqrt{2}}&\dfrac{(1+\gamma_5)\gamma_1}{\sqrt{2}}&\dfrac{(1+\gamma_5)\gamma_2}{\sqrt{2}}&\dfrac{(1+\gamma_5)\gamma_3}{\sqrt{2}}\\
-\dfrac{(1-\gamma_5)\gamma_0}{\sqrt{2}}&-\dfrac{(1+\gamma_5)\gamma_0}{\sqrt{2}}&0&\gamma_0\gamma_1&\gamma_0\gamma_2&\gamma_0\gamma_3\\
-\dfrac{(1-\gamma_5)\gamma_1}{\sqrt{2}}&-\dfrac{(1+\gamma_5)\gamma_1}{\sqrt{2}}&\gamma_1\gamma_0&0&\gamma_1\gamma_2&\gamma_1\gamma_3\\
-\dfrac{(1-\gamma_5)\gamma_2}{\sqrt{2}}&-\dfrac{(1+\gamma_5)\gamma_2}{\sqrt{2}}&\gamma_2\gamma_0&\gamma_2\gamma_1&0&\gamma_2\gamma_3\\
-\dfrac{(1-\gamma_5)\gamma_3}{\sqrt{2}}&-\dfrac{(1+\gamma_5)\gamma_3}{\sqrt{2}}&\gamma_3\gamma_0&\gamma_3\gamma_1&\gamma_3\gamma_2&0
\end{pmatrix}.
\label{eq:ifd-dirac-jab}
\end{align}
\endgroup
Thus $\rho_\gamma(J_{-2,-1})=-i\gamma_5/2=-D$ and the
translation and special-conformal entries reproduce
$J_{-1,\mu}=P_\mu/\sqrt{2}$ and
$J_{-2,\mu}=-\mathfrak{K}_\mu/\sqrt{2}$, respectively. The Clifford
algebra consequently reproduces the IFD commutators in
Eq.~\eqref{algebrasimp} and Table~\ref{tabelIFD} \cite{Dirac1936,Hepner1962}.
}
\section{Interpolating conformal algebra in \texorpdfstring{$(3+1)$}{Lg} dimensions}
\label{conformalsimpler}
The IFD projective algebra is interpolated by the following $6\times6$
transformation matrix:
\begin{align}
    (\mathcal{R}_{\hat{a}}^{b})_{6\times6}=(\mathcal{R}_{\hat{a}}^{b})^T_{6\times6}=\begin{pmatrix}
    1&0&0&0&0&0\\
    0&1&0&0&0&0\\
    0&0&\cos{\delta}&0&0&\sin{\delta}\\
    0&0&0&1&0&0\\
    0&0&0&0&1&0\\
    0&0&\sin{\delta}&0&0&-\cos{\delta}
    \end{pmatrix}
\end{align}
Then in interpolation form $J_{\hat{a}\hat{b}}$ becomes $J_{\hat{a}\hat{b}}=\mathcal{R}_{\hat{a}}^{~{c}}J_{cd}\mathcal{R}_{~\hat{b}}^{{d}}$, that is
\begin{align}
    J_{\hat{a}\hat{b}}&=\begin{pmatrix}
    0&-D&-\frac{\mathfrak{K}^{\hat{+}}}{\sqrt{2}}&-\frac{\mathfrak{K}_1}{\sqrt{2}}&-\frac{\mathfrak{K}_2}{\sqrt{2}}&-\frac{\mathfrak{K}^{\hat{-}}}{\sqrt{2}}\\
    D&0&\frac{P_{\hat{+}}}{\sqrt{2}}&\frac{P_{1}}{\sqrt{2}}&\frac{P_{2}}{\sqrt{2}}&\frac{P_{\hat{-}}}{\sqrt{2}}\\
    \frac{\mathfrak{K}^{\hat{+}}}{\sqrt{2}}&-\frac{P_{\hat{+}}}{\sqrt{2}}&0 & {\mathcal{D}}^{\itP{1}} & {\mathcal{D}}^{\itP{2}} & -K_3\\
    \frac{\mathfrak{K}_1}{\sqrt{2}}&-\frac{P_{1}}{\sqrt{2}}&-{\mathcal{D}}^{\itP{1}} & 0 & {J}^{3} & -{\mathcal{K}}^{\itP{1}}\\
    \frac{\mathfrak{K}_2}{\sqrt{2}}&-\frac{P_{2}}{\sqrt{2}}&-{\mathcal{D}}^{\itP{2}} & -{J}^{3} & 0 & -{\mathcal{K}}^{\itP{2}}\\
    \frac{\mathfrak{K}^{\hat{-}}}{\sqrt{2}}&-\frac{P_{\hat{-}}}{\sqrt{2}}&K_3 & {\mathcal{K}}^{\itP{1}} & {\mathcal{K}}^{\itP{2}} & 0
  \end{pmatrix}_{(6\times6)}\label{Jhat+hat-}
\end{align}
Then the simplified conformal algebra in interpolation is:
  \begin{align}
      \left[J_{{\hat{a}}{\hat{b}}},J_{{\hat{c}}{\hat{d}}}\right]=-i\left(g_{{\hat{b}}{\hat{d}}}J_{{\hat{a}}{\hat{c}}}-g_{{\hat{b}}{\hat{c}}}J_{{\hat{a}}{\hat{d}}}+g_{{\hat{a}}{\hat{c}}}J_{{\hat{b}}{\hat{d}}}-g_{{\hat{a}}{\hat{d}}}J_{{\hat{b}}{\hat{c}}}\right)\label{simplesrint}
  \end{align}
where, 
\begin{align}
    g_{\hat{a},\hat{b}}&=\begin{pmatrix}
  0&-1&0&0&0&0\\
  -1&0&0&0&0&0\\
  0&0&\mathbb{C}&0&0&\mathbb{S}\\
  0&0&0&-1&0&0\\
  0&0&0&0&-1&0\\
  0&0&\mathbb{S}&0&0&-\mathbb{C}\\
  \end{pmatrix}_{6\times6}\label{metricghat}
\end{align}
The algebra Eq.\eqref{simplesrint} with the above interpolating $6\times6$ metric will reproduce 105 explicit commutation relations in the interpolation between IFD and LFD mentioned in Table~\ref{tabelinterpolation}.

The interpolation in projective spacetime naturally produces superscript special conformal generators  ($\mathfrak{K}^{\hat{\pm}}$), which are related to the subscript special conformal generators as $\mathfrak{K}^{\hat{\pm}}=g^{\hat{\pm}\hat{\pm}}\mathfrak{K}_{\hat{\pm}}+g^{\hat{\pm}\hat{\mp}}\mathfrak{K}_{\hat{\mp}}$, where the metric is defined in Eq.\eqref{eqn:g_munu_interpolation}, explicitly,
\begin{align}
    \begin{pmatrix}
        \mathfrak{K}^{\hat{+}}\\
        \mathfrak{K}^{\hat{-}}
    \end{pmatrix}=&\begin{pmatrix}
        \mathbb{C}&\mathbb{S}\\
        \mathbb{S}&-\mathbb{C}
    \end{pmatrix}\begin{pmatrix}
        \mathfrak{K}_{\hat{+}}\\
        \mathfrak{K}_{\hat{-}}
    \end{pmatrix}.
\end{align}

\begin{table}[htpb!]
\centering
\renewcommand{\arraystretch}{2} % Adds vertical padding for fractions
\caption{The 15 conformal generators in the interpolation form alongside their Instant-Form Dynamics (IFD) and Light-Front Dynamics (LFD) limits.}
\label{tab:generator_limits}
\begin{tabular}{@{}ccc@{}}
\toprule
\makecell{\textbf{Interpolating} \\ \textbf{generators}} & 
\makecell{\textbf{IFD limit} \\ ($\delta \to 0$)} & 
\makecell{\textbf{LFD limit} \\ ($\delta \to \frac{\pi}{4}$)} \\
\midrule

% Translations P
$\begin{aligned}
P_{\hat{+}} &= P_0\cos\delta + P_3\sin\delta \\
P_{\hat{-}} &= P_0\sin\delta - P_3\cos\delta
\end{aligned}$ &
$\begin{aligned}
P_{\hat{+}} &\to P_0 \\
P_{\hat{-}} &\to -P_3
\end{aligned}$ &
$\begin{aligned}
P_{\hat{+}} &\to P_{+} = \frac{P_0 + P_3}{\sqrt{2}} \\
P_{\hat{-}} &\to P_{-} = \frac{P_0 - P_3}{\sqrt{2}}
\end{aligned}$ \\

\midrule

% Special Conformal K
$\begin{aligned}
\mathfrak{K}_{\hat{+}} &= \mathfrak{K}_0\cos\delta - \mathfrak{K}_3\sin\delta \\
\mathfrak{K}_{\hat{-}} &= \mathfrak{K}_0\sin\delta + \mathfrak{K}_3\cos\delta
\end{aligned}$ &
$\begin{aligned}
\mathfrak{K}_{\hat{+}} &\to \mathfrak{K}_0 \\
\mathfrak{K}_{\hat{-}} &\to \mathfrak{K}_3
\end{aligned}$ &
$\begin{aligned}
\mathfrak{K}_{\hat{+}} &\to \mathfrak{K}_{+} = \frac{\mathfrak{K}_0 - \mathfrak{K}_3}{\sqrt{2}} \\
\mathfrak{K}_{\hat{-}} &\to \mathfrak{K}_{-} = \frac{\mathfrak{K}_0 + \mathfrak{K}_3}{\sqrt{2}}
\end{aligned}$ \\

\midrule

% Dilatations & Longitudinal Boosts
$\begin{aligned}
D_{\hat{+}} &= D\cos\delta + K_3\sin\delta \\
D_{\hat{-}} &= D\sin\delta - K_3\cos\delta
\end{aligned}$ &
$\begin{aligned}
D_{\hat{+}} &\to D \\
D_{\hat{-}} &\to -K_3
\end{aligned}$ &
$\begin{aligned}
D_{\hat{+}} &\to D_{+} = \frac{D + K_3}{\sqrt{2}} \\
D_{\hat{-}} &\to D_{-} = \frac{D - K_3}{\sqrt{2}}
\end{aligned}$ \\

\midrule

% Transverse K (Mathcal K)
$\begin{aligned}
\mathcal{K}^{\hat{1}} &= -K^1\sin\delta - J^2\cos\delta \\
\mathcal{K}^{\hat{2}} &= J^1\cos\delta - K^2\sin\delta
\end{aligned}$ &
$\begin{aligned}
\mathcal{K}^{\hat{1}} &\to -J^2 \\
\mathcal{K}^{\hat{2}} &\to J^1
\end{aligned}$ &
$\begin{aligned}
\mathcal{K}^{\hat{1}} &\to -E^1 = \frac{-K^1 - J^2}{\sqrt{2}} \\
\mathcal{K}^{\hat{2}} &\to -E^2 = \frac{J^1 - K^2}{\sqrt{2}}
\end{aligned}$ \\

\midrule

% Transverse D (Mathcal D)
$\begin{aligned}
\mathcal{D}^{\hat{1}} &= -K^1\cos\delta + J^2\sin\delta \\
\mathcal{D}^{\hat{2}} &= -J^1\sin\delta - K^2\cos\delta
\end{aligned}$ &
$\begin{aligned}
\mathcal{D}^{\hat{1}} &\to -K^1 \\
\mathcal{D}^{\hat{2}} &\to -K^2
\end{aligned}$ &
$\begin{aligned}
\mathcal{D}^{\hat{1}} &\to -F^1 = \frac{-K^1 + J^2}{\sqrt{2}} \\
\mathcal{D}^{\hat{2}} &\to -F^2 = \frac{-J^1 - K^2}{\sqrt{2}}
\end{aligned}$ \\

\midrule

% Rotations and Transverse invariant generators (To total 15)
$J^3 = J^3$ & $J^3 \to J^3$ & $J^3 \to J^3$ \\
$P_{1,2} = P_{1,2}$ & $P_{1,2} \to P_{1,2}$ & $P_{1,2} \to P_{1,2}$ \\
$\mathfrak{K}_{1,2} = \mathfrak{K}_{1,2}$ & $\mathfrak{K}_{1,2} \to \mathfrak{K}_{1,2}$ & $\mathfrak{K}_{1,2} \to \mathfrak{K}_{1,2}$ \\

\bottomrule
\end{tabular}
\end{table}



\startlandscape
\begingroup
\fontsize{8}{9}\selectfont
\setlength{\tabcolsep}{1pt}
\renewcommand{\arraystretch}{2.5}
\noindent
\begin{minipage}[t]{\textwidth}
\centering
\captionof{table}{Full $15 \times 15$ conformal algebra in the interpolation form}
\label{tab:full_interpolation}
\resizebox{\textwidth}{!}{%
\begin{tabular}{ |c||c|c|c|c|c|c|c|c|c|c|c|c|c|c|c| }
\hline
\rule{0pt}{16pt}  & $\mathfrak{K}_{\hat{+}}$ & $P_{\hat{+}}$ & $D_{\hat{+}}$ & $\mathfrak{K}_{\hat{-}}$ & $P_{\hat{-}}$ & $D_{\hat{-}}$ & $P_{1}$ & $P_{2}$ & $\mathfrak{K}_{1}$ & $\mathfrak{K}_{2}$ & $\mathcal{D}^{\hat{1}}$ & $\mathcal{D}^{\hat{2}}$ & $\mathcal{K}^{\hat{1}}$ & $\mathcal{K}^{\hat{2}}$ & $J^{3}$ \\
\hline
\hline
\rule{0pt}{16pt} $\mathfrak{K}_{\hat{+}}$ & $0$ & $-2i[(c+\mathbb{S}s)D_{\hat{+}}+(s-\mathbb{S}c)D_{\hat{-}}]$ & $-i[(c+\mathbb{S}s)\mathfrak{K}_{\hat{+}}+(s-\mathbb{S}c)\mathfrak{K}_{\hat{-}}]$ & $0$ & $2i\mathbb{C}[sD_{\hat{+}}-cD_{\hat{-}}]$ & $i\mathbb{C}[s\mathfrak{K}_{\hat{+}}-c\mathfrak{K}_{\hat{-}}]$ & $-2i[\mathbb{S}\mathcal{K}^{\hat{1}} + \mathbb{C}\mathcal{D}^{\hat{1}}]$ & $-2i[\mathbb{S}\mathcal{K}^{\hat{2}} + \mathbb{C}\mathcal{D}^{\hat{2}}]$ & $0$ & $0$ & $i\mathfrak{K}_{1}$ & $i\mathfrak{K}_{2}$ & $0$ & $0$ & $0$ \\
\hline
\rule{0pt}{16pt} $P_{\hat{+}}$ & $2i[(c+\mathbb{S}s)D_{\hat{+}}+(s-\mathbb{S}c)D_{\hat{-}}]$ & $0$ & $i[(c+\mathbb{S}s)P_{\hat{+}}+(s-\mathbb{S}c)P_{\hat{-}}]$ & $-2i\mathbb{C}[sD_{\hat{+}}-cD_{\hat{-}}]$ & $0$ & $-i\mathbb{C}[sP_{\hat{+}}-cP_{\hat{-}}]$ & $0$ & $0$ & $-2i[\mathbb{S}\mathcal{K}^{\hat{1}} + \mathbb{C}\mathcal{D}^{\hat{1}}]$ & $-2i[\mathbb{S}\mathcal{K}^{\hat{2}} + \mathbb{C}\mathcal{D}^{\hat{2}}]$ & $i\mathbb{C} P_{1}$ & $i\mathbb{C} P_{2}$ & $i\mathbb{S} P_{1}$ & $i\mathbb{S} P_{2}$ & $0$ \\
\hline
\rule{0pt}{16pt} $D_{\hat{+}}$ & $i[(c+\mathbb{S}s)\mathfrak{K}_{\hat{+}}+(s-\mathbb{S}c)\mathfrak{K}_{\hat{-}}]$ & $-i[(c+\mathbb{S}s)P_{\hat{+}}+(s-\mathbb{S}c)P_{\hat{-}}]$ & $0$ & $-i\mathbb{C}[s\mathfrak{K}_{\hat{+}}-c\mathfrak{K}_{\hat{-}}]$ & $i\mathbb{C}[sP_{\hat{+}}-cP_{\hat{-}}]$ & $0$ & $-ic P_{1}$ & $-ic P_{2}$ & $ic\mathfrak{K}_{1}$ & $ic\mathfrak{K}_{2}$ & $is[\mathbb{C}\mathcal{K}^{\hat{1}} - \mathbb{S}\mathcal{D}^{\hat{1}}]$ & $is[\mathbb{C}\mathcal{K}^{\hat{2}} - \mathbb{S}\mathcal{D}^{\hat{2}}]$ & $is[\mathbb{S}\mathcal{K}^{\hat{1}} + \mathbb{C}\mathcal{D}^{\hat{1}}]$ & $is[\mathbb{S}\mathcal{K}^{\hat{2}} + \mathbb{C}\mathcal{D}^{\hat{2}}]$ & $0$ \\
\hline
\rule{0pt}{16pt} $\mathfrak{K}_{\hat{-}}$ & $0$ & $2i\mathbb{C}[sD_{\hat{+}}-cD_{\hat{-}}]$ & $i\mathbb{C}[s\mathfrak{K}_{\hat{+}}-c\mathfrak{K}_{\hat{-}}]$ & $0$ & $-2i[(c-\mathbb{S}s)D_{\hat{+}}+(s+\mathbb{S}c)D_{\hat{-}}]$ & $-i[(c-\mathbb{S}s)\mathfrak{K}_{\hat{+}}+(s+\mathbb{S}c)\mathfrak{K}_{\hat{-}}]$ & $2i[\mathbb{C}\mathcal{K}^{\hat{1}} - \mathbb{S}\mathcal{D}^{\hat{1}}]$ & $2i[\mathbb{C}\mathcal{K}^{\hat{2}} - \mathbb{S}\mathcal{D}^{\hat{2}}]$ & $0$ & $0$ & $0$ & $0$ & $i\mathfrak{K}_{1}$ & $i\mathfrak{K}_{2}$ & $0$ \\
\hline
\rule{0pt}{16pt} $P_{\hat{-}}$ & $-2i\mathbb{C}[sD_{\hat{+}}-cD_{\hat{-}}]$ & $0$ & $-i\mathbb{C}[sP_{\hat{+}}-cP_{\hat{-}}]$ & $2i[(c-\mathbb{S}s)D_{\hat{+}}+(s+\mathbb{S}c)D_{\hat{-}}]$ & $0$ & $i[(c-\mathbb{S}s)P_{\hat{+}}+(s+\mathbb{S}c)P_{\hat{-}}]$ & $0$ & $0$ & $2i[\mathbb{C}\mathcal{K}^{\hat{1}} - \mathbb{S}\mathcal{D}^{\hat{1}}]$ & $2i[\mathbb{C}\mathcal{K}^{\hat{2}} - \mathbb{S}\mathcal{D}^{\hat{2}}]$ & $i\mathbb{S} P_{1}$ & $i\mathbb{S} P_{2}$ & $-i\mathbb{C} P_{1}$ & $-i\mathbb{C} P_{2}$ & $0$ \\
\hline
\rule{0pt}{16pt} $D_{\hat{-}}$ & $-i\mathbb{C}[s\mathfrak{K}_{\hat{+}}-c\mathfrak{K}_{\hat{-}}]$ & $i\mathbb{C}[sP_{\hat{+}}-cP_{\hat{-}}]$ & $0$ & $i[(c-\mathbb{S}s)\mathfrak{K}_{\hat{+}}+(s+\mathbb{S}c)\mathfrak{K}_{\hat{-}}]$ & $-i[(c-\mathbb{S}s)P_{\hat{+}}+(s+\mathbb{S}c)P_{\hat{-}}]$ & $0$ & $-is P_{1}$ & $-is P_{2}$ & $is\mathfrak{K}_{1}$ & $is\mathfrak{K}_{2}$ & $-ic[\mathbb{C}\mathcal{K}^{\hat{1}} - \mathbb{S}\mathcal{D}^{\hat{1}}]$ & $-ic[\mathbb{C}\mathcal{K}^{\hat{2}} - \mathbb{S}\mathcal{D}^{\hat{2}}]$ & $-ic[\mathbb{S}\mathcal{K}^{\hat{1}} + \mathbb{C}\mathcal{D}^{\hat{1}}]$ & $-ic[\mathbb{S}\mathcal{K}^{\hat{2}} + \mathbb{C}\mathcal{D}^{\hat{2}}]$ & $0$ \\
\hline
\rule{0pt}{16pt} $P_{1}$ & $2i[\mathbb{S}\mathcal{K}^{\hat{1}} + \mathbb{C}\mathcal{D}^{\hat{1}}]$ & $0$ & $ic P_{1}$ & $-2i[\mathbb{C}\mathcal{K}^{\hat{1}} - \mathbb{S}\mathcal{D}^{\hat{1}}]$ & $0$ & $is P_{1}$ & $0$ & $0$ & $-2i[cD_{\hat{+}} + sD_{\hat{-}}]$ & $-2iJ^{3}$ & $i P_{\hat{+}}$ & $0$ & $i P_{\hat{-}}$ & $0$ & $-i P_{2}$ \\
\hline
\rule{0pt}{16pt} $P_{2}$ & $2i[\mathbb{S}\mathcal{K}^{\hat{2}} + \mathbb{C}\mathcal{D}^{\hat{2}}]$ & $0$ & $ic P_{2}$ & $-2i[\mathbb{C}\mathcal{K}^{\hat{2}} - \mathbb{S}\mathcal{D}^{\hat{2}}]$ & $0$ & $is P_{2}$ & $0$ & $0$ & $2iJ^{3}$ & $-2i[cD_{\hat{+}} + sD_{\hat{-}}]$ & $0$ & $i P_{\hat{+}}$ & $0$ & $i P_{\hat{-}}$ & $i P_{1}$ \\
\hline
\rule{0pt}{16pt} $\mathfrak{K}_{1}$ & $0$ & $2i[\mathbb{S}\mathcal{K}^{\hat{1}} + \mathbb{C}\mathcal{D}^{\hat{1}}]$ & $-ic\mathfrak{K}_{1}$ & $0$ & $-2i[\mathbb{C}\mathcal{K}^{\hat{1}} - \mathbb{S}\mathcal{D}^{\hat{1}}]$ & $-is\mathfrak{K}_{1}$ & $2i[cD_{\hat{+}} + sD_{\hat{-}}]$ & $-2iJ^{3}$ & $0$ & $0$ & $i\mathfrak{K}_{\hat{+}}$ & $0$ & $i\mathfrak{K}_{\hat{-}}$ & $0$ & $-i\mathfrak{K}_{2}$ \\
\hline
\rule{0pt}{16pt} $\mathfrak{K}_{2}$ & $0$ & $2i[\mathbb{S}\mathcal{K}^{\hat{2}} + \mathbb{C}\mathcal{D}^{\hat{2}}]$ & $-ic\mathfrak{K}_{2}$ & $0$ & $-2i[\mathbb{C}\mathcal{K}^{\hat{2}} - \mathbb{S}\mathcal{D}^{\hat{2}}]$ & $-is\mathfrak{K}_{2}$ & $2iJ^{3}$ & $2i[cD_{\hat{+}} + sD_{\hat{-}}]$ & $0$ & $0$ & $0$ & $i\mathfrak{K}_{\hat{+}}$ & $0$ & $i\mathfrak{K}_{\hat{-}}$ & $i\mathfrak{K}_{1}$ \\
\hline
\rule{0pt}{16pt} $\mathcal{D}^{\hat{1}}$ & $-i\mathfrak{K}_{1}$ & $-i\mathbb{C} P_{1}$ & $-is[\mathbb{C}\mathcal{K}^{\hat{1}} - \mathbb{S}\mathcal{D}^{\hat{1}}]$ & $0$ & $-i\mathbb{S} P_{1}$ & $ic[\mathbb{C}\mathcal{K}^{\hat{1}} - \mathbb{S}\mathcal{D}^{\hat{1}}]$ & $-i P_{\hat{+}}$ & $0$ & $-i\mathfrak{K}_{\hat{+}}$ & $0$ & $0$ & $-i\mathbb{C} J^{3}$ & $-i[sD_{\hat{+}} - cD_{\hat{-}}]$ & $-i\mathbb{S} J^{3}$ & $-i\mathcal{D}^{\hat{2}}$ \\
\hline
\rule{0pt}{16pt} $\mathcal{D}^{\hat{2}}$ & $-i\mathfrak{K}_{2}$ & $-i\mathbb{C} P_{2}$ & $-is[\mathbb{C}\mathcal{K}^{\hat{2}} - \mathbb{S}\mathcal{D}^{\hat{2}}]$ & $0$ & $-i\mathbb{S} P_{2}$ & $ic[\mathbb{C}\mathcal{K}^{\hat{2}} - \mathbb{S}\mathcal{D}^{\hat{2}}]$ & $0$ & $-i P_{\hat{+}}$ & $0$ & $-i\mathfrak{K}_{\hat{+}}$ & $i\mathbb{C} J^{3}$ & $0$ & $i\mathbb{S} J^{3}$ & $-i[sD_{\hat{+}} - cD_{\hat{-}}]$ & $i\mathcal{D}^{\hat{1}}$ \\
\hline
\rule{0pt}{16pt} $\mathcal{K}^{\hat{1}}$ & $0$ & $-i\mathbb{S} P_{1}$ & $-is[\mathbb{S}\mathcal{K}^{\hat{1}} + \mathbb{C}\mathcal{D}^{\hat{1}}]$ & $-i\mathfrak{K}_{1}$ & $i\mathbb{C} P_{1}$ & $ic[\mathbb{S}\mathcal{K}^{\hat{1}} + \mathbb{C}\mathcal{D}^{\hat{1}}]$ & $-i P_{\hat{-}}$ & $0$ & $-i\mathfrak{K}_{\hat{-}}$ & $0$ & $i[sD_{\hat{+}} - cD_{\hat{-}}]$ & $-i\mathbb{S} J^{3}$ & $0$ & $i\mathbb{C} J^{3}$ & $-i\mathcal{K}^{\hat{2}}$ \\
\hline
\rule{0pt}{16pt} $\mathcal{K}^{\hat{2}}$ & $0$ & $-i\mathbb{S} P_{2}$ & $-is[\mathbb{S}\mathcal{K}^{\hat{2}} + \mathbb{C}\mathcal{D}^{\hat{2}}]$ & $-i\mathfrak{K}_{2}$ & $i\mathbb{C} P_{2}$ & $ic[\mathbb{S}\mathcal{K}^{\hat{2}} + \mathbb{C}\mathcal{D}^{\hat{2}}]$ & $0$ & $-i P_{\hat{-}}$ & $0$ & $-i\mathfrak{K}_{\hat{-}}$ & $i\mathbb{S} J^{3}$ & $i[sD_{\hat{+}} - cD_{\hat{-}}]$ & $-i\mathbb{C} J^{3}$ & $0$ & $i\mathcal{K}^{\hat{1}}$ \\
\hline
\rule{0pt}{16pt} $J^{3}$ & $0$ & $0$ & $0$ & $0$ & $0$ & $0$ & $i P_{2}$ & $-i P_{1}$ & $i\mathfrak{K}_{2}$ & $-i\mathfrak{K}_{1}$ & $i\mathcal{D}^{\hat{2}}$ & $-i\mathcal{D}^{\hat{1}}$ & $i\mathcal{K}^{\hat{2}}$ & $-i\mathcal{K}^{\hat{1}}$ & $0$ \\
\hline
\end{tabular}}
\label{tabelinterpolation}
\end{minipage}

\newpage

\noindent
\begin{minipage}[t]{\textwidth}
\centering
\renewcommand{\arraystretch}{1.5}
\adjustbox{max width=\textwidth, max totalheight=0.85\textheight}{%
\begin{tabular}{ |c||c|c|c|c|c|c|c|c|c|c|c|c|c|c|c| }
\hline
\rule{0pt}{16pt} & $\mathfrak{K}_{+}$ & $P_{+}$ & $D_{+}$ & $\mathfrak{K}_{-}$ & $P_{-}$ & $D_{-}$ & $P_{1}$ & $P_{2}$ & $\mathfrak{K}_{1}$ & $\mathfrak{K}_{2}$ & $F^{1}$ & $F^{2}$ & $E^{1}$ & $E^{2}$ & $J^{3}$ \\
\hline
\hline
\rule{0pt}{16pt} $\mathfrak{K}_{+}$ & $0$ & $-2i\sqrt{2}D_{+}$ & $-i\sqrt{2}\mathfrak{K}_{+}$ & $0$ & $0$ & $0$ & $2iE^{1}$ & $2iE^{2}$ & $0$ & $0$ & $-i\mathfrak{K}_{1}$ & $-i\mathfrak{K}_{2}$ & $0$ & $0$ & $0$ \\
\hline
\rule{0pt}{16pt} $P_{+}$ & $2i\sqrt{2}D_{+}$ & $0$ & $i\sqrt{2}P_{+}$ & $0$ & $0$ & $0$ & $0$ & $0$ & $2iE^{1}$ & $2iE^{2}$ & $0$ & $0$ & $-iP_{1}$ & $-iP_{2}$ & $0$ \\
\hline
\rule{0pt}{16pt} $D_{+}$ & $i\sqrt{2}\mathfrak{K}_{+}$ & $-i\sqrt{2}P_{+}$ & $0$ & $0$ & $0$ & $0$ & $-i\frac{1}{\sqrt{2}}P_{1}$ & $-i\frac{1}{\sqrt{2}}P_{2}$ & $i\frac{1}{\sqrt{2}}\mathfrak{K}_{1}$ & $i\frac{1}{\sqrt{2}}\mathfrak{K}_{2}$ & $-i\frac{1}{\sqrt{2}}F^{1}$ & $-i\frac{1}{\sqrt{2}}F^{2}$ & $i\frac{1}{\sqrt{2}}E^{1}$ & $i\frac{1}{\sqrt{2}}E^{2}$ & $0$ \\
\hline
\rule{0pt}{16pt} $\mathfrak{K}_{-}$ & $0$ & $0$ & $0$ & $0$ & $-2i\sqrt{2}D_{-}$ & $-i\sqrt{2}\mathfrak{K}_{-}$ & $2iF^{1}$ & $2iF^{2}$ & $0$ & $0$ & $0$ & $0$ & $-i\mathfrak{K}_{1}$ & $-i\mathfrak{K}_{2}$ & $0$ \\
\hline
\rule{0pt}{16pt} $P_{-}$ & $0$ & $0$ & $0$ & $2i\sqrt{2}D_{-}$ & $0$ & $i\sqrt{2}P_{-}$ & $0$ & $0$ & $2iF^{1}$ & $2iF^{2}$ & $-iP_{1}$ & $-iP_{2}$ & $0$ & $0$ & $0$ \\
\hline
\rule{0pt}{16pt} $D_{-}$ & $0$ & $0$ & $0$ & $i\sqrt{2}\mathfrak{K}_{-}$ & $-i\sqrt{2}P_{-}$ & $0$ & $-i\frac{1}{\sqrt{2}}P_{1}$ & $-i\frac{1}{\sqrt{2}}P_{2}$ & $i\frac{1}{\sqrt{2}}\mathfrak{K}_{1}$ & $i\frac{1}{\sqrt{2}}\mathfrak{K}_{2}$ & $i\frac{1}{\sqrt{2}}F^{1}$ & $i\frac{1}{\sqrt{2}}F^{2}$ & $-i\frac{1}{\sqrt{2}}E^{1}$ & $-i\frac{1}{\sqrt{2}}E^{2}$ & $0$ \\
\hline
\rule{0pt}{16pt} $P_{1}$ & $-2iE^{1}$ & $0$ & $i\frac{1}{\sqrt{2}}P_{1}$ & $-2iF^{1}$ & $0$ & $i\frac{1}{\sqrt{2}}P_{1}$ & $0$ & $0$ & $-i\sqrt{2}(D_{+}+D_{-})$ & $-2iJ^{3}$ & $-iP_{+}$ & $0$ & $-iP_{-}$ & $0$ & $-iP_{2}$ \\
\hline
\rule{0pt}{16pt} $P_{2}$ & $-2iE^{2}$ & $0$ & $i\frac{1}{\sqrt{2}}P_{2}$ & $-2iF^{2}$ & $0$ & $i\frac{1}{\sqrt{2}}P_{2}$ & $0$ & $0$ & $2iJ^{3}$ & $-i\sqrt{2}(D_{+}+D_{-})$ & $0$ & $-iP_{+}$ & $0$ & $-iP_{-}$ & $iP_{1}$ \\
\hline
\rule{0pt}{16pt} $\mathfrak{K}_{1}$ & $0$ & $-2iE^{1}$ & $-i\frac{1}{\sqrt{2}}\mathfrak{K}_{1}$ & $0$ & $-2iF^{1}$ & $-i\frac{1}{\sqrt{2}}\mathfrak{K}_{1}$ & $i\sqrt{2}(D_{+}+D_{-})$ & $-2iJ^{3}$ & $0$ & $0$ & $-i\mathfrak{K}_{+}$ & $0$ & $-i\mathfrak{K}_{-}$ & $0$ & $-i\mathfrak{K}_{2}$ \\
\hline
\rule{0pt}{16pt} $\mathfrak{K}_{2}$ & $0$ & $-2iE^{2}$ & $-i\frac{1}{\sqrt{2}}\mathfrak{K}_{2}$ & $0$ & $-2iF^{2}$ & $-i\frac{1}{\sqrt{2}}\mathfrak{K}_{2}$ & $2iJ^{3}$ & $i\sqrt{2}(D_{+}+D_{-})$ & $0$ & $0$ & $0$ & $-i\mathfrak{K}_{+}$ & $0$ & $-i\mathfrak{K}_{-}$ & $i\mathfrak{K}_{1}$ \\
\hline
\rule{0pt}{16pt} $F^{1}$ & $i\mathfrak{K}_{1}$ & $0$ & $i\frac{1}{\sqrt{2}}F^{1}$ & $0$ & $iP_{1}$ & $-i\frac{1}{\sqrt{2}}F^{1}$ & $iP_{+}$ & $0$ & $i\mathfrak{K}_{+}$ & $0$ & $0$ & $0$ & $i\frac{1}{\sqrt{2}}(D_{+}-D_{-})$ & $-iJ^{3}$ & $-iF^{2}$ \\
\hline
\rule{0pt}{16pt} $F^{2}$ & $i\mathfrak{K}_{2}$ & $0$ & $i\frac{1}{\sqrt{2}}F^{2}$ & $0$ & $iP_{2}$ & $-i\frac{1}{\sqrt{2}}F^{2}$ & $0$ & $iP_{+}$ & $0$ & $i\mathfrak{K}_{+}$ & $0$ & $0$ & $iJ^{3}$ & $i\frac{1}{\sqrt{2}}(D_{+}-D_{-})$ & $iF^{1}$ \\
\hline
\rule{0pt}{16pt} $E^{1}$ & $0$ & $iP_{1}$ & $-i\frac{1}{\sqrt{2}}E^{1}$ & $i\mathfrak{K}_{1}$ & $0$ & $i\frac{1}{\sqrt{2}}E^{1}$ & $iP_{-}$ & $0$ & $i\mathfrak{K}_{-}$ & $0$ & $-i\frac{1}{\sqrt{2}}(D_{+}-D_{-})$ & $-iJ^{3}$ & $0$ & $0$ & $-iE^{2}$ \\
\hline
\rule{0pt}{16pt} $E^{2}$ & $0$ & $iP_{2}$ & $-i\frac{1}{\sqrt{2}}E^{2}$ & $i\mathfrak{K}_{2}$ & $0$ & $i\frac{1}{\sqrt{2}}E^{2}$ & $0$ & $iP_{-}$ & $0$ & $i\mathfrak{K}_{-}$ & $iJ^{3}$ & $-i\frac{1}{\sqrt{2}}(D_{+}-D_{-})$ & $0$ & $0$ & $iE^{1}$ \\
\hline
\rule{0pt}{16pt} $J^{3}$ & $0$ & $0$ & $0$ & $0$ & $0$ & $0$ & $iP_{2}$ & $-iP_{1}$ & $i\mathfrak{K}_{2}$ & $-i\mathfrak{K}_{1}$ & $iF^{2}$ & $-iF^{1}$ & $iE^{2}$ & $-iE^{1}$ & $0$ \\
\hline
\end{tabular}}
\captionof{table}{Conformal algebra $(3+1)$ in the LFD limit ($\delta \to \pi/4$)}
\label{tabelLFD}
\end{minipage}
\endgroup
\finishlandscape


where $\mathbb{S}=\sin2\delta$,  $\mathbb{C}=\cos2\delta$, $s=\sin{\delta}$ and $c=\cos{\delta}$. With the IFD limit $\mathbb{C}\rightarrow1~\&~\mathbb{S}\rightarrow0$ and LFD limit $\mathbb{C}\rightarrow0~\&~\mathbb{S}\rightarrow1$, the algebra Eq.\eqref{simplesrint} will reproduce the explicit commutation relations presented in Table~\ref{tabelIFD} and Table~\ref{tabelLFD}, respectively.


\section{\texorpdfstring{$6\times6$}{TEXT} Projective Spacetime Representations}
\label{sec:projectivespace}
To characterize the individual properties of the 15 conformal generators, namely, which generators are kinematic or dynamic, we may correspond the projective spacetime representations 
given by Eq.(\ref{Jab}) with
the interpolating $(3+1)$ dimensional conformal algebra
discussed in the previous sections, especially Section~\ref{conformalsimpler}. 
Kinematic generators leave the time defined in the given form of the dynamics invariant.
Consequently, the individual time-ordered amplitudes in the respective time-ordered processes are invariant under the transformations provided by kinematic generators \cite{Ji2001}.
Determining which generators are kinematic in each form of dynamics is useful, as maximizing the number of kinematic generators correspondingly saves dynamical effort when solving relativistic quantum field theories, such as QED and QCD~ \cite{ji2023relativistic}. As discussed in Section~\ref{conformalsimpler}, 
LFD saves dynamical effort in solving relativistic quantum field theoretic problems due to the character change of the longitudinal boost operator $K_3$ from dynamical in IFD to kinematic in LFD, as exemplified in 't Hooft model computations  \cite{THOOFT1974461, Ji2021QCD, ji2023relativistic}.

 The conformal algebra given by Eq.\eqref{simplesrint} in the interpolating form of projective spacetime implies that $J_{\hat{a}\hat{b}}$ can be written as
\begin{align}
    J_{\hat{a}\hat{b}}=i(X_{\hat{a}}\partial_{\hat{b}}-X_{\hat{b}}\partial_{\hat{a}}),
\end{align}
where $\hat{a},\hat{b}\in\{-2,-1,\hat{+},1,2,\hat{-}\}$, $\partial_{\hat{a}}\equiv\frac{\partial}{\partial X^{\hat{a}}}$, and $X_{\hat{a}}$ is a six-dimensional projective embedding vector for the physical $(3+1)$-dimensional spacetime. With the condition that its lightcone is preserved under the transformations generated by $J_{\hat{a}\hat{b}}$, we write the infinitesimal conformal transformations as
\begin{align}
R^{\hat{a}}_{~\hat{b}}=g^{\hat{a}}_{~\hat{b}}-\frac{i}{2}\omega^{\hat{c}\hat{d}}(J_{\hat{c}\hat{d}})^{\hat{a}}_{~\hat{b}},\label{R_conformalTransformations}
\end{align}
where $\omega^{\hat{a}\hat{b}}=-\omega^{\hat{b}\hat{a}}$. Here, we may regard $R^{\hat{a}}_{~\hat{b}}$ as the expansion of $e^{-\frac{i}{2}\omega^{\hat{c}\hat{d}}(J_{\hat{c}\hat{d}})^{\hat{a}}_{~\hat{b}}}$ up to the first order in $\omega$, as can be seen in the expansion given by 
$e^{-\frac{i}{2}\omega^{\hat{c}\hat{d}}(J_{\hat{c}\hat{d}})^{\hat{a}}_{~\hat{b}}} \approx  g^{\hat{a}}_{~\hat{b}}-\frac{i}{2}\omega^{\hat{c}\hat{d}}(J_{\hat{c}\hat{d}})^{\hat{a}}_{~\hat{b}}+\mathcal{O}(\omega^2)$.
The generator representation $(J_{\hat{c}\hat{d}})^{\hat{a}}_{~\hat{b}}$ can be obtained by
\begin{align}
    (J_{\hat{c}\hat{d}})^{\hat{a}}_{~\hat{b}}=(J_{\hat{c}\hat{d}})^{\hat{a}\hat{f}}g_{\hat{f}\hat{b}},
\end{align}
where $(J_{{\hat{c}}{\hat{d}}})^{{\hat{a}}{\hat{b}}}=i(g_{\hat{c}}^{\hat{a}} g_{\hat{d}}^{\hat{b}}-g_{\hat{c}}^{\hat{b}} g_{\hat{d}}^{\hat{a}})$. The representation matrices of the conformal generators are defined by taking the first index as a superscript and the second as a subscript:

\begin{eqnarray}
    &&(D_{\hat{+}})^{\hat{a}}_{~\hat{b}}\equiv -\cos{\delta}(J_{-2-1})^{\hat{a}}_{~\hat{b}}-\sin{\delta}(J_{\hat{+}\hat{-}})^{\hat{a}}_{~\hat{b}}~;\nonumber\\
    &&(D_{\hat{-}})^{\hat{a}}_{~\hat{b}}\equiv -\sin{\delta}(J_{-2-1})^{\hat{a}}_{~\hat{b}}+\cos{\delta}(J_{\hat{+}\hat{-}})^{\hat{a}}_{~\hat{b}}~;\nonumber\\
    &&\frac{-(\mathfrak{K}^{\hat{\mu}})^{\hat{a}}_{~\hat{b}}}{\sqrt{2}}\equiv (J_{-2\hat{\mu}})^{\hat{a}}_{~\hat{b}}~;
    \frac{(P_{\hat{\mu}})^{\hat{a}}_{~\hat{b}}}{\sqrt{2}}\equiv (J_{-1\hat{\mu}})^{\hat{a}}_{~\hat{b}}~.
    \label{implicit6x6}
\end{eqnarray}

The explicit $6\times6$ matrix representations are then given by
\begin{align}
    D_{\hat{+}}&=\begin{pmatrix}
        -i\cos{(\delta)}&0&0&0&0&0\\
        0&i\cos{(\delta)}&0&0&0&0\\
        0&0&i\mathbb{S}\sin{(\delta)}&0&0&-i\mathbb{C}\sin{(\delta)}\\
        0&0&0&0&0&0\\
        0&0&0&0&0&0\\
        0&0&-i\mathbb{C}\sin{(\delta)}&0&0&-i\mathbb{S}\sin{(\delta)}
    \end{pmatrix};\nonumber\\
    D_{\hat{-}}&=\begin{pmatrix}
        -i\sin{(\delta)}&0&0&0&0&0\\
        0&i\sin{(\delta)}&0&0&0&0\\
        0&0&-i\mathbb{S}\cos{(\delta)}&0&0&i\mathbb{C}\cos{(\delta)}\\
        0&0&0&0&0&0\\
        0&0&0&0&0&0\\
        0&0&i\mathbb{C}\cos{(\delta)}&0&0&i\mathbb{S}\cos{(\delta)}
    \end{pmatrix};\nonumber
\end{align}
\begin{align}
    \mathfrak{K}^{\hat{+}}&=\sqrt{2}\begin{pmatrix}
        0&0&0&0&0&0\\
        0&0&i&0&0&0\\
        i\mathbb{C}&0&0&0&0&0\\
        0&0&0&0&0&0\\
        0&0&0&0&0&0\\
        i\mathbb{S}&0&0&0&0&0
    \end{pmatrix};
    ~\mathfrak{K}^{\hat{-}}=\sqrt{2}\begin{pmatrix}
        0&0&0&0&0&0\\
        0&0&0&0&0&i\\
        i\mathbb{S}&0&0&0&0&0\\
        0&0&0&0&0&0\\
        0&0&0&0&0&0\\
        -i\mathbb{C}&0&0&0&0&0
    \end{pmatrix};\nonumber
    \end{align}
    \begin{align}
     P_{\hat{+}}&=\sqrt{2}\begin{pmatrix}
        0&0&-i&0&0&0\\
        0&0&0&0&0&0\\
        0&-i\mathbb{C}&0&0&0&0\\
        0&0&0&0&0&0\\
        0&0&0&0&0&0\\
        0&-i\mathbb{S}&0&0&0&0
    \end{pmatrix}~;P_{\hat{-}}=\sqrt{2}\begin{pmatrix}
        0&0&0&0&0&-i\\
        0&0&0&0&0&0\\
        0&-i\mathbb{S}&0&0&0&0\\
        0&0&0&0&0&0\\
        0&0&0&0&0&0\\
        0&i\mathbb{C}&0&0&0&0
    \end{pmatrix}~;\nonumber
    \end{align}
    \begin{align}
        \mathcal{K}^{1}=\begin{pmatrix}
        0&0&0&0&0&0\\
        0&0&0&0&0&0\\
        0&0&0&i\mathbb{S}&0&0\\
        0&0&0&0&0&i\\
        0&0&0&0&0&0\\
        0&0&0&-i\mathbb{C}&0&0
    \end{pmatrix}~;~~~~~~\mathcal{K}^{2}=\begin{pmatrix}
        0&0&0&0&0&0\\
        0&0&0&0&0&0\\
        0&0&0&0&i\mathbb{S}&0\\
        0&0&0&0&0&0\\
        0&0&0&0&0&i\\
        0&0&0&0&-i\mathbb{C}&0
    \end{pmatrix}~;\nonumber
    \end{align}
    \begin{align}
    \mathcal{D}^{1}&=\begin{pmatrix}
        0&0&0&0&0&0\\
        0&0&0&0&0&0\\
        0&0&0&i\mathbb{C}&0&0\\
        0&0&i&0&0&0\\
        0&0&0&0&0&0\\
        0&0&0&i\mathbb{S}&0&0
    \end{pmatrix}~;~~~~~~\mathcal{D}^{2}=\begin{pmatrix}
        0&0&0&0&0&0\\
        0&0&0&0&0&0\\
        0&0&0&0&i\mathbb{C}&0\\
        0&0&0&0&0&0\\
        0&0&i&0&0&0\\
        0&0&0&0&i\mathbb{S}&0
    \end{pmatrix}~;\nonumber
    \end{align}
    \begin{align}
        \mathfrak{K}_1=\sqrt{2}\begin{pmatrix}
    0&0&0&0&0&0\\
    0&0&0&i&0&0\\
    0&0&0&0&0&0\\
    -i&0&0&0&0&0\\
    0&0&0&0&0&0\\
    0&0&0&0&0&0
    \end{pmatrix}~;
    \mathfrak{K}_2&=\sqrt{2}\begin{pmatrix}
    0&0&0&0&0&0\\
    0&0&0&0&i&0\\
    0&0&0&0&0&0\\
    0&0&0&0&0&0\\
    -i&0&0&0&0&0\\
    0&0&0&0&0&0
    \end{pmatrix}~;\nonumber
    \end{align}
    \begin{align}
        P_1=\sqrt{2}\begin{pmatrix}
    0&0&0&-i&0&0\\
    0&0&0&0&0&0\\
    0&0&0&0&0&0\\
    0&i&0&0&0&0\\
    0&0&0&0&0&0\\
    0&0&0&0&0&0
    \end{pmatrix}~;~~~~~~
    P_2=\sqrt{2}\begin{pmatrix}
    0&0&0&0&-i&0\\
    0&0&0&0&0&0\\
    0&0&0&0&0&0\\
    0&0&0&0&0&0\\
    0&i&0&0&0&0\\
    0&0&0&0&0&0
    \end{pmatrix}~.\nonumber
    \end{align}
    \begin{align}
        J^3&=\begin{pmatrix}
    0&0&0&0&0&0\\
    0&0&0&0&0&0\\
    0&0&0&0&0&0\\
  0&0&0&0 & -i & 0\\
  0&0&0&i & 0 & 0\\
  0&0&0&0 & 0 & 0
\end{pmatrix}\label{matrix-rep-conformal-generators}
    \end{align}
%

These explicit projective spacetime matrices satisfy the full conformal algebra Eq.\eqref{simplesrint} in interpolation. 

Operating the exponentiated conformal transformations $e^{-\frac{i}{2}\omega^{\hat{c}\hat{d}}(J_{\hat{c}\hat{d}})^{\hat{a}}_{~\hat{b}}}$ on $X_{\hat{b}}$, we obtain the transformed six-dimensional projective embedding vector $X^{\prime}_{\hat{a}}$ as
\begin{align}
\label{exponential-conformal-tranf}
  X^{\prime}_{\hat{a}}&=e^{-\frac{i}{2}\omega^{\hat{c}\hat{d}}(J_{\hat{c}\hat{d}})^{\hat{b}}_{~\hat{a}}}X_{\hat{b}},
\end{align}
and the inner product
\begin{align}
{X^{\prime}}_{\hat{a}}{X^{\prime}}^{\hat{a}} &= 
e^{-\frac{i}{2}\omega^{\hat{c}\hat{d}}(J_{\hat{c}\hat{d}})^{\hat{b}}_{~\hat{a}}}X_{\hat{b}}e^{-\frac{i}{2}\omega^{\hat{c}\hat{d}}(J_{\hat{c}\hat{d}})^{\hat{a}}_{~\hat{b}}}X^{\hat{b}}\nonumber\\
&= e^{-\frac{i}{2}\omega^{\hat{c}\hat{d}}(J_{\hat{c}\hat{d}})^{\hat{b}}_{~\hat{a}}}e^{-\frac{i}{2}\omega^{\hat{c}\hat{d}}(J_{\hat{c}\hat{d}})^{\hat{a}}_{~\hat{b}}}X_{\hat{b}}X^{\hat{b}}\nonumber\\
&= e^{-\frac{i}{2}\omega^{\hat{c}\hat{d}}((J_{\hat{c}\hat{d}})^{\hat{b}}_{~\hat{a}}+(J_{\hat{c}\hat{d}})^{\hat{a}}_{~\hat{b}})}X_{\hat{b}}X^{\hat{b}}\nonumber\\
&= X_{\hat{b}}X^{\hat{b}},
\end{align}
where we used $(J_{{\hat{c}}{\hat{d}}})^{\hat{a}}_{~\hat{b}}=i(g_{\hat{c}}^{\hat{a}} g_{{\hat{d}}{\hat{b}}}-g_{{\hat{c}}{\hat{b}}} g_{\hat{d}}^{\hat{a}})$ to obtain $(J_{\hat{c}\hat{d}})^{\hat{a}}_{~\hat{b}}=-(J_{\hat{c}\hat{d}})^{\hat{b}}_{~\hat{a}}$.
%$(J_{\hat{c}\hat{d}})^{\hat{b}}_{~\hat{a}}+(J_{\hat{c}\hat{d}})^{\hat{a}}_{~\hat{b}}=0$. 
This shows that the conformal transformations, even in the interpolating form, preserve the angles and the ratios of lengths of two vectors before and after the transformations. Now, we assign the following six-dimensional projective embedding vector:
\begin{align}
    {X}_{-1}&=\frac{-\lambda}{\sqrt{2}}, \nonumber\\
    {X}_{-2}&=\frac{-\lambda}{\sqrt{2}}(x^{\hat{\mu}}x_{\hat{\mu}}), \nonumber \\
     {X}_{\hat{\mu}}&=\lambda x_{\hat{\mu}},\label{X-mu}
\end{align}
provides the lightcone of the projective spacetime, i.e., ${X}_{\hat{a}}{X}^{\hat{a}}=0$, as explicitly shown below: 
\begin{align}
    {X}_{\hat{a}}{X}^{\hat{a}}&={X}_{-2}{X}^{-2}+{X}_{-1}{X}^{-1}+{X}_{\hat{\mu}}{X}^{\hat{\mu}} \nonumber \\
    &=-2{X}_{-2}{X}_{-1}+{X}^{\hat{\mu}}{X}_{\hat{\mu}}\nonumber \\
     &=-2\frac{-\lambda}{\sqrt{2}}(x^{\hat{\mu}}x_{\hat{\mu}})\frac{-\lambda}{\sqrt{2}}+\lambda^2{x}^{\hat{\mu}}{x}_{\hat{\mu}}\nonumber \\
     &=-\lambda^2{x}^{\hat{\mu}}{x}_{\hat{\mu}}+\lambda^2{x}^{\hat{\mu}}{x}_{\hat{\mu}}\nonumber \\
     &=0,
\end{align}
where ${X}^{-1}=-{X}_{-2}$, ${X}^{-2}=-{X}_{-1}$ and $X^{\hat \mu}=g^{\hat{\mu}\hat{\nu}}X_{\hat{\nu}}$ due to the metric
given by Eq.(\ref{metricghat}).
From Eq.\eqref{X-mu}, we can then find 
the dictionary connecting the six-dimensional projective embedding vector $X_{\hat{a}}$ and the physical $(3+1)$-dimensional spacetime $x_{\hat{\mu}}$, which reads
\begin{align}
    x_{\hat{\mu}}&=-\frac{1}{\sqrt{2}}\frac{{X}_{\hat{\mu}}}{ {X}_{-1}}.
\end{align}
Corresponding to $x_{\hat{\mu}}$, we have  
\begin{align}
    x^{\hat{\mu}}&=-\frac{1}{\sqrt{2}}\frac{{X}^{\hat{\mu}}}{ {X}_{-1}}.
\end{align}
Then, the inverse spacetime of $x^{\hat{\mu}}$ is given by $\frac{x_{\hat{\mu}}}{x^2}=-\frac{1}{\sqrt{2}}\frac{{X}_{\hat{\mu}}}{{X}_{-2}}$, with $x^2=\frac{{X}_{-2}}{{X}_{-1}}$. Under the transformation given by Eq.(\ref{exponential-conformal-tranf}), we compute the transformation of interpolating time given by
\begin{align}
    x^{\hat{+}\prime}&=-\frac{1}{\sqrt{2}}\frac{{X}^{\hat{+}\prime}}{ {X}^{\prime}_{-1}},
\end{align}
for each and every interpolating conformal operation according to Eq.(\ref{exponential-conformal-tranf}). We summarized the results of our computation in Table~\ref{tablexprime}.
From these results, we can determine which generators leave the interpolating time intact. We call a generator kinematic when $x^{\hat{+}\prime}=\lambda x^{\hat{+}}$, where $\lambda$ is independent of the spacetime coordinates (and may depend on the group parameter). Transformations with coordinate-dependent multiplicative factors are therefore classified as dynamic. For $0\leq \delta < \pi/4$, the seven kinematic generators are $\mathcal{K}^{\hat{1}}$, $\mathcal{K}^{\hat{2}}$, $J^{3}$, $P_{1}$, $P_{2}$, $P_{\hat -}=\sqrt{2}J_{-1\hat{-}}$, and $D=\left(D_{\hat{+}}c+D_{\hat{-}}s\right)=J_{-1 -2}$.
The first six of these generators leave $x^{\hat{+}}$ exactly unchanged. The scale factor for the dilatation $D$ is $e^{-\alpha}$, with $\alpha=\omega^{-1-2}$ in Eq.(\ref{exponential-conformal-tranf}), as follows straightforwardly from the $6\times6$ matrix representations of $D_{\hat{+}}$ and $D_{\hat{-}}$ in Eq.(\ref{matrix-rep-conformal-generators}).
Taking the respective limit of the interpolating parameter $\delta$, namely,  
$\delta \to 0$ and $\delta \to \pi/4$, we find the corresponding number of kinematic generators in the IFD and the LFD, respectively. 
%
\begin{table}[h!]
        %\centering
        \resizebox{\textwidth}{!}{
        \begin{tabular}{|c||c|c|}
        \hline
             \rule{0pt}{16pt} Generators & $x^{\hat{+}\prime}$\\
             \hline
            \rule{0pt}{16pt} $\mathfrak{K}^{\hat{+}}$ & $x^{\hat{+}\prime}=\frac{x^{\hat{+}}-b^{\hat{+}}(x)^2}{1-2\mathbb{C}b^{\hat{+}}x^{\hat{+}}-2\mathbb{S}b^{\hat{+}}x^{\hat{-}}+\mathbb{C}(b^{\hat{+}})^2(x)^2}$\\
             \hline
            \rule{0pt}{16pt} $\mathfrak{K}^{\hat{-}}$ & $x^{\hat{+}\prime}=\frac{x^{\hat{+}}}{1+2\mathbb{C}b^{\hat{-}}x^{\hat{-}}-2\mathbb{S}b^{\hat{-}}x^{\hat{+}}-\mathbb{C}(b^{\hat{-}})^2(x)^2}$ \\
             \hline
             \rule{0pt}{16pt} $P_{\hat{+}}$& $x^{\hat{+}\prime}=x^{\hat{+}}+a^{\hat{+}} $\\
             \hline
             \rule{0pt}{16pt} $P_{\hat{-}}$& $x^{\hat{+}\prime}=x^{\hat{+}}$  \\
             \hline
             \rule{0pt}{16pt} $D_{\hat{+}}$ & $x^{\hat{+}\prime} = e^{-\alpha_{\hat{+}}\cos{(\delta)}}\left[\left(\cosh[\alpha_{\hat{+}}\sin{(\delta)}]+\mathbb{S} \sinh[\alpha_{\hat{+}}\sin{(\delta)}]\right)x^{\hat{+}} -\mathbb{C} \sinh[\alpha_{\hat{+}}\sin{(\delta)}]x^{\hat{-}}\right]$ \\
             \hline
             \rule{0pt}{16pt} $D_{\hat{-}}$ & $x^{\hat{+}\prime} = e^{-\alpha_{\hat{-}}\sin{(\delta)}}\left[\left(\cosh[\alpha_{\hat{-}}\cos{(\delta)}]-\mathbb{S}\sinh[\alpha_{\hat{-}}\cos{(\delta)}]\right)x^{\hat{+}}+\mathbb{C}\sinh[\alpha_{\hat{-}}\cos{(\delta)}] x^{\hat{-}}\right]$ \\
             \hline
             \rule{0pt}{16pt} $J^{3}$ & $ x^{\hat{+}\prime}=x^{\hat{+}}$ \\
             \hline
             \rule{0pt}{16pt} $\mathfrak{K}_{1}$ & $x^{\hat{+}\prime}=\frac{x^{{\hat{+}}}}{1-2b^{1}x_{1}-(b^{1})^{2}(x)^2}$ \\
             \hline
             \rule{0pt}{16pt} $\mathfrak{K}_{2}$ &  $x^{\hat{+}\prime}=\frac{x^{{\hat{+}}}}{1-2b^{2}x_{2}-(b^{2})^{2}(x)^2}$ \\
             \hline
              \rule{0pt}{16pt} $P_{1}$ & $ x^{\hat{+}\prime}=x^{\hat{+}}$  \\
              \hline
              \rule{0pt}{16pt} $P_{2}$ & $ x^{\hat{+}\prime}=x^{\hat{+}}$  \\
              \hline
               \rule{0pt}{16pt} $\mathcal{D}^{1}$ &  $ x^{\hat{+}\prime}=\cosh{(\sqrt{\mathbb{C}}\eta^{1})}x^{\hat{+}}+\frac{\mathbb{S}}{\mathbb{C}}\left(\cosh{(\sqrt{\mathbb{C}}\eta^{1})}-1\right)x^{\hat{-}}+\frac{1}{\sqrt{\mathbb{C}}}\sinh{(\sqrt{\mathbb{C}}\eta^{1})}x_{{1}}$ \\
               \hline
               \rule{0pt}{16pt} $\mathcal{D}^{2}$ & $ x^{\hat{+}\prime}=\cosh{(\sqrt{\mathbb{C}}\eta^{2})}x^{\hat{+}}+\frac{\mathbb{S}}{\mathbb{C}}\left(\cosh{(\sqrt{\mathbb{C}}\eta^{2})}-1\right)x^{\hat{-}}+\frac{1}{\sqrt{\mathbb{C}}}\sinh{(\sqrt{\mathbb{C}}\eta^{2})}x_{{2}}$ \\
               \hline
               \rule{0pt}{16pt} $\mathcal{K}^{1}$ & $ x^{\hat{+}\prime}=x^{\hat{+}}$ \\
               \hline
               \rule{0pt}{16pt} $\mathcal{K}^{2}$ & $ x^{\hat{+}\prime}=x^{\hat{+}}$ \\
               \hline
        \end{tabular}}
        \caption{Transformation of interpolating time under each conformal generator in $(3+1)$ dimensions. Each row denotes a one-parameter transformation $X'=e^{-i\lambda_G G}X$, with the parameter associated with the displayed generator $G$: $b^{\hat{\pm}}$, $b^1$, and $b^2$ for special conformal transformations; $a^{\hat{\pm}}$, $a^1$, and $a^2$ for translations; $\alpha_{\hat{\pm}}$ for $D_{\hat{\pm}}$; $\eta^1$ and $\eta^2$ for $\mathcal{D}^{1}$ and $\mathcal{D}^{2}$; $\phi^1$ and $\phi^2$ for $\mathcal{K}^{1}$ and $\mathcal{K}^{2}$; and $\theta$ for $J^3$. The interpolating parameter between IFD and LFD is $\delta$.}
        \label{tablexprime}
    \end{table}
\FloatBarrier

\newcommand{\IFDTimeTable}{%
    \begin{table}[!htbp]
      \centering
      \small
        \begin{tabular}{|c||c|}
        \hline
             \rule{0pt}{16pt} Generators & $x^{{0}\prime}$\\
             \hline
             \hline
            \rule{0pt}{16pt} $\mathfrak{K}_{{0}}$ & $x^{{0}\prime}=\frac{x^{{0}}-b^{{0}}(x)^2}{1-2b^{{0}}x^{{0}}+(b^{{0}})^2(x)^2}$\\
             \hline
            \rule{0pt}{16pt} $-\mathfrak{K}_{{3}}$ & $x^{{0}\prime}=\frac{x^{{0}}}{1+2b^{{3}}x^{{3}}+(b^{{3}})^2(x)^2}$ \\
             \hline
             \rule{0pt}{16pt} $P_{{0}}$& $x^{{0}\prime}=x^{{0}}+a^{{0}} $\\
             \hline
             \rule{0pt}{16pt} $-P_{{3}}$& $x^{{0}\prime}=x^{{0}}$  \\
             \hline
             \rule{0pt}{16pt} $D$ & $x^{{0}\prime} = e^{-\alpha}x^{0}$ \\
             \hline
             \rule{0pt}{16pt} $-K_{3}$ & $x^{{0}\prime} = \cosh[\alpha]x^{0}-\sinh[\alpha] x^{3}$ \\
             \hline
             \rule{0pt}{16pt} $J^{3}$ & $ x^{{0}\prime}=x^{{0}}$ \\
             \hline
             \rule{0pt}{16pt} $\mathfrak{K}_{1}$ & $x^{{0}\prime}=\frac{x^{{{0}}}}{1-2b^{1}x_{1}-(b^{1})^{2}(x)^2}$ \\
             \hline
             \rule{0pt}{16pt} $\mathfrak{K}_{2}$ &  $x^{{0}\prime}=\frac{x^{{{0}}}}{1-2b^{2}x_{2}-(b^{2})^{2}(x)^2}$ \\
             \hline
              \rule{0pt}{16pt} $P_{1}$ & $ x^{{0}\prime}=x^{{0}}$  \\
              \hline
              \rule{0pt}{16pt} $P_{2}$ & $ x^{{0}\prime}=x^{{0}}$  \\
              \hline
               \rule{0pt}{16pt} $-K^{1}$ &  $ x^{{0}\prime}=\cosh{(\eta^{1})}x^{{0}}+\sinh{(\eta^{1})}x_{{1}}$ \\
               \hline
               \rule{0pt}{16pt} $-K^{2}$ & $ x^{{0}\prime}=\cosh{(\eta^{2})}x^{{0}}+\sinh{(\eta^{2})}x_{{2}}$ \\
               \hline
               \rule{0pt}{16pt} $-J^{2}$ & $ x^{{0}\prime}=x^{{0}}$ \\
               \hline
               \rule{0pt}{16pt} $J^{1}$ & $ x^{{0}\prime}=x^{{0}}$ \\
               \hline
        \end{tabular}
        \caption{Transformation of the IFD time in $(3+1)$d}
        \label{tablex0prime}
    \end{table}%
}

\newcommand{\LFDTimeTable}{%
    \begin{table}[!htbp]
        \centering
        \small
        \begin{tabular}{|c||c|}
        \hline
             \rule{0pt}{16pt} Generators & $x^{{+}\prime}$\\
             \hline
             \hline
            \rule{0pt}{16pt} $\mathfrak{K}_{{-}}$ & $x^{{+}\prime}=\frac{x^{{+}}-b^{{+}}(2x^{+}x^{-}-\vec{x}^{\perp}\vec{x}^{\perp})}{1-2b^{{+}}x^{{-}}}$\\
             \hline
            \rule{0pt}{16pt} $\mathfrak{K}_{{+}}$ & $x^{{+}\prime}=\frac{x^{{+}}}{1-2b^{{-}}x^{{+}}}$ \\
             \hline
             \rule{0pt}{16pt} $P_{{+}}$& $x^{{+}\prime}=x^{{+}}+a^{{+}} $\\
             \hline
             \rule{0pt}{16pt} $P_{{-}}$& $x^{{+}\prime}=x^{{+}}$  \\
             \hline
             \rule{0pt}{16pt} $D_{{+}}$ & $x^{{+}\prime} = e^{-\sqrt{2}\alpha}x^{{+}}$ \\
             \hline
             \rule{0pt}{16pt} $D_{{-}}$ & $x^{{+}\prime} = x^{{+}}$ \\
             \hline
             \rule{0pt}{16pt} $J^{3}$ & $ x^{{+}\prime}=x^{{+}}$ \\
             \hline
             \rule{0pt}{16pt} $\mathfrak{K}_{1}$ & $x^{{+}\prime}=\frac{x^{{{+}}}}{1-2b^{1}x_{1}-(b^{1})^{2}(x)^2}$ \\
             \hline
             \rule{0pt}{16pt} $\mathfrak{K}_{2}$ &  $x^{{+}\prime}=\frac{x^{{{+}}}}{1-2b^{2}x_{2}-(b^{2})^{2}(x)^2}$ \\
             \hline
              \rule{0pt}{16pt} $P_{1}$ & $ x^{{+}\prime}=x^{{+}}$  \\
              \hline
              \rule{0pt}{16pt} $P_{2}$ & $ x^{{+}\prime}=x^{{+}}$  \\
              \hline
               \rule{0pt}{16pt} $-F^{1}$ &  $x^{{+}\prime}=x^{{+}}+\frac{(\eta^{1})^2}{2}x^{{-}}+\eta^{1}x_{{1}}$ \\
               \hline
               \rule{0pt}{16pt} $-F^{2}$ & $ x^{{+}\prime}=x^{{+}}+\frac{(\eta^{2})^2}{2}x^{{-}}+\eta^{2}x_{{2}}$ \\
               \hline
               \rule{0pt}{16pt} $-E^{1}$ & $ x^{{+}\prime}=x^{{+}}$ \\
               \hline
               \rule{0pt}{16pt}$-E^{2}$ & $ x^{{+}\prime}=x^{{+}}$ \\
               \hline
        \end{tabular}
        \caption{Transformation of LFD time in $(3+1)$d}
        \label{tablexplusprime}
    \end{table}%
}

\newcommand{\KinematicGeneratorTable}{%
    \begin{table}[!htbp]
    \resizebox{\textwidth}{!}{%
    \begin{tabular}{lcc}
	Interpolation angle & Kinematic & Dynamic \\
	\hline
	\rule{0pt}{3ex} $\delta=0$ & $\mathcal{K}^{\hat{1}}=-J^{2},~ \mathcal{K}^{\hat{2}}=J^{1},~ J^{3}, P^{1}, P^{2}, P^{3}$, $D$ & $\mathcal{D}^{\hat{1}}=-K^{1},~ \mathcal{D}^{\hat{2}}=-K^{2},~ K^{3}, P^{0}$, $\mathfrak{K}_{{0}}$, $\mathfrak{K}_{{1}}$, $\mathfrak{K}_{{2}}$, $\mathfrak{K}_{{3}}$\\
	$0<\delta<\pi/4$ &  $\mathcal{K}^{\hat{1}}, \mathcal{K}^{\hat{2}}, J^{3}, P^{1}, P^{2}$, $P_{\mT}$, $D=\left(D_{\hat{+}}c+D_{\hat{-}}s\right)$ &  $\mathcal{D}^{\hat{1}}, \mathcal{D}^{\hat{2}}$, $K_3=\left(D_{\hat{+}}s-D_{\hat{-}}c\right)$,  $P_{\pT}$, $\mathfrak{K}_{\hat{+}}$, $\mathfrak{K}_{1}$, $\mathfrak{K}_{2}$, $\mathfrak{K}_{\hat{-}}$\\
	$\delta=\pi/4$ & $\mathcal{K}^{\hat{1}}=-E^{1},~ \mathcal{K}^{\hat{2}}=-E^{2},~ J^{3}, P^{1}, P^{2}$, $P_{-}$, $D_{+}$, $D_{-}$& $\mathcal{D}^{\hat{1}}=-F^{1},~ \mathcal{D}^{\hat{2}}=-F^{2},~ P_{+}$, $\mathfrak{K}_{{+}}$, $\mathfrak{K}_{{-}}$, $\mathfrak{K}_{{1}}$, $\mathfrak{K}_{{2}}$\\
	%\hline
	 %\hline
      \end{tabular}%
    }
    \caption{\label{tab:Kinematic_and_dynamic_generators_for_different_interoplation_angles_conformal1+1}Kinematic and dynamic conformal generators for different interpolation angles in $(3+1)$ dimensions. Here, kinematic means that the interpolating time is unchanged up to a scale factor independent of the spacetime coordinates.}
\end{table}%
}
%


In the limit to $\delta=0$, i.e., the IFD ($\mathbb{C}=1;~\mathbb{S}=0$), we find the results shown in Table~\ref{tablex0prime} for the transformed $x^{0\prime}$ 
for every $(3+1)$d conformal generator. Consistent with Table~\ref{tablexprime}, the seven kinematic generators are $\mathcal{K}^{\hat{1}}=-J^{2}$, $\mathcal{K}^{\hat{2}}=J^{1}$, $J^{3}$, $P^{1}$, $P^{2}$, $P^{3}=-P_{3}$, and $D$.

\IFDTimeTable
\FloatBarrier

In the limit to $\delta=\pi/4$, i.e., the LFD ($\mathbb{C}=0;~\mathbb{S}=1$), we find the results shown in Table~\ref{tablexplusprime} for the transformed $x^{+\prime}$ for each $(3+1)$d conformal generator. The eight kinematic generators are $\mathcal{K}^{\hat{1}}=-E^{1}$, $\mathcal{K}^{\hat{2}}=-E^{2}$, $J^{3}$, $P^{1}$, $P^{2}$, $P_{-}=\frac{P_{0}-P_{3}}{\sqrt{2}}$, $D_{+}=\frac{(D+K_3)}{\sqrt{2}}$, and $D_{-}=\frac{(D-K_3)}{\sqrt{2}}$. Equivalently, the latter two generators span the kinematic directions $D$ and $K_3$, where $K_3=\left(D_{\hat{+}}s-D_{\hat{-}}c\right)=J_{\hat{-}\hat{+}}$ at $\delta=\pi/4$.

\LFDTimeTable
\FloatBarrier

Finally, we contrast the behavior of the special conformal transformations (SCT) in $(3+1)$ dimensions with the $(1+1)$ dimensional case. In $(1+1)$ dimensions, the transverse coordinates vanish ($\vec{x}^{\perp} = 0$), and the transformation under $\mathfrak{K}_{-}$ simplifies to $x^{+\prime} = x^{+}$, making it an additional kinematic generator. However, in $(3+1)$ dimensions, the presence of the transverse term $\vec{x}^{\perp}\cdot\vec{x}^{\perp}$ in the numerator of $x^{+\prime}$ (as shown in Table~\ref{tab:Kinematic_and_dynamic_generators_for_different_interoplation_angles_conformal1+1}) means that the light-front time is no longer invariant under $\mathfrak{K}_{-}$. Consequently, $\mathfrak{K}_{-}$ transforms from being kinematic in $(1+1)$d to becoming dynamical in $(3+1)$d.

\KinematicGeneratorTable
\FloatBarrier

As a result, when moving from IFD to LFD in $(3+1)$ dimensions, the kinematic subalgebra gains the longitudinal-boost direction $K_3$. As summarized in Table~\ref{tab:Kinematic_and_dynamic_generators_for_different_interoplation_angles_conformal1+1}, out of the 15 conformal generators in $(3+1)$ dimensions, 8 generators are kinematic in LFD, compared to 7 in IFD. This demonstrates that LFD still maximizes the total number of kinematic generators in $(3+1)$ dimensions, albeit with a different dynamical balance than in lower dimensions.


\section{Summary and Conclusion}
\label{summary-conclusion}

In this work, we presented the interpolating conformal algebra between IFD and LFD in $(3+1)$ dimensions. We elaborated on the $(3+1)$ dimensional conformal algebra using the interpolating projective spacetime representation and showed that LFD maximizes the number of kinematic generators. As discussed in this work, identifying which conformal generators become kinematic in LFD is useful for understanding the utility of LFD in hadron phenomenology, thereby saving considerable dynamical effort when solving QCD.

Building on the correspondence between $(1+1)$d quantum field theory and AdS$_3$, we plan to analyze the holographic dual correspondence between $(3+1)$d QCD and AdS$_5$ by further interpolating between IFD and LFD. With this algebraic framework, we can evaluate light-front QCD, which maximizes the number of kinematic generators and links first-principle QCD to the light-front quark model for hadronic phenomenology. Handling relativistic interactions among hadron constituents is relevant for understanding the advantages of different forms of relativistic dynamics~\cite{Ji2025}.


\acknowledgments{This work was supported by the U.S. Department of Energy (Grant No. DE-FG02-03ER41260).
The National Energy Research Scientific Computing Center (NERSC) supported by the Office of Science of the U.S. Department of Energy under Contract No. DE-AC02-05CH11231 is also acknowledged. 
}
\appendix
\setcounter{section}{0}
\section{Derivation of the Interpolating-Time Transformations}
\label{app:interpolating-time-derivations}

This appendix gives the essential derivations of the fifteen entries in
Table~\ref{tablexprime}. For each individual transformation, we apply the
corresponding matrix in Eq.~\eqref{matrix-rep-conformal-generators} directly
to the interpolating projective space-time vector. The resulting interpolating
projective space-time components are substituted
in the defining ratios for $x_{\hat{+}}'$ and $x_{\hat{-}}'$, and then
$x^{\hat{+}\prime}$ is obtained by raising the index as in
Section~\ref{sec:projectivespace}.
Only one generator parameter is nonzero in each subsection.

\subsection{Longitudinal special conformal transformations}

\subsubsection{\texorpdfstring{$\mathfrak{K}^{\hat{+}}$}{K hat plus}}

The transformation of the interpolating projective space-time vector by
$\mathfrak{K}^{\hat{+}}$ is
\begin{align}
\begin{pmatrix}
X_{-2}'\\ X_{-1}'\\ X_{\hat{+}}'\\ X_1'\\ X_2'\\ X_{\hat{-}}'
\end{pmatrix}
&=e^{-ib^{\hat{+}}\mathfrak{K}^{\hat{+}}}
\begin{pmatrix}
X_{-2}\\ X_{-1}\\ X_{\hat{+}}\\ X_1\\ X_2\\ X_{\hat{-}}
\end{pmatrix}
=\begin{pmatrix}
1 & 0 & 0 & 0 & 0 & 0\\
\mathbb{C} (b^{\hat{+}})^{2} & 1 & \sqrt{2} b^{\hat{+}} & 0 & 0 & 0\\
\sqrt{2} \mathbb{C} b^{\hat{+}} & 0 & 1 & 0 & 0 & 0\\
0 & 0 & 0 & 1 & 0 & 0\\
0 & 0 & 0 & 0 & 1 & 0\\
\sqrt{2} \mathbb{S} b^{\hat{+}} & 0 & 0 & 0 & 0 & 1
\end{pmatrix}
\begin{pmatrix}
X_{-2}\\ X_{-1}\\ X_{\hat{+}}\\ X_1\\ X_2\\ X_{\hat{-}}
\end{pmatrix}\\
&=\begin{pmatrix}
X_{-2}\\
X_{-1}+\sqrt{2}b^{\hat{+}}X_{\hat{+}}
+\mathbb{C}(b^{\hat{+}})^2X_{-2}\\
X_{\hat{+}}+\sqrt{2}\mathbb{C}b^{\hat{+}}X_{-2}\\
X_1\\
X_2\\
X_{\hat{-}}+\sqrt{2}\mathbb{S}b^{\hat{+}}X_{-2}
\end{pmatrix}.
\end{align}
\begin{align}
x_{\hat{+}}'
&=-\frac{1}{\sqrt{2}}\frac{X_{\hat{+}}'}{X_{-1}'}
=-\frac{1}{\sqrt{2}}
\frac{X_{\hat{+}}/X_{-1}+\sqrt{2}\mathbb{C}b^{\hat{+}}X_{-2}/X_{-1}}
{1+\sqrt{2}b^{\hat{+}}X_{\hat{+}}/X_{-1}
 +\mathbb{C}(b^{\hat{+}})^2X_{-2}/X_{-1}}\nonumber\\
&=\frac{x_{\hat{+}}-\mathbb{C}b^{\hat{+}}x^2}
{1-2b^{\hat{+}}x_{\hat{+}}+\mathbb{C}(b^{\hat{+}})^2x^2}, \nonumber\\
x_{\hat{-}}'
&=-\frac{1}{\sqrt{2}}\frac{X_{\hat{-}}'}{X_{-1}'}
=-\frac{1}{\sqrt{2}}
\frac{X_{\hat{-}}/X_{-1}+\sqrt{2}\mathbb{S}b^{\hat{+}}X_{-2}/X_{-1}}
{1+\sqrt{2}b^{\hat{+}}X_{\hat{+}}/X_{-1}
 +\mathbb{C}(b^{\hat{+}})^2X_{-2}/X_{-1}}\nonumber\\
&=\frac{x_{\hat{-}}-\mathbb{S}b^{\hat{+}}x^2}
{1-2b^{\hat{+}}x_{\hat{+}}+\mathbb{C}(b^{\hat{+}})^2x^2}.
\end{align}
Consequently,
\begin{align}
x^{\hat{+}\prime}
&=\mathbb{C}x_{\hat{+}}'+\mathbb{S}x_{\hat{-}}'\nonumber\\
&=\frac{x^{\hat{+}}-b^{\hat{+}}x^2}
{1-2\mathbb{C}b^{\hat{+}}x^{\hat{+}}
 -2\mathbb{S}b^{\hat{+}}x^{\hat{-}}
 +\mathbb{C}(b^{\hat{+}})^2x^2},
\end{align}
which is the first entry of Table~\ref{tablexprime}.

\subsubsection{\texorpdfstring{$\mathfrak{K}^{\hat{-}}$}{K hat minus}}

The transformation of the interpolating projective space-time vector by
$\mathfrak{K}^{\hat{-}}$ is
\begin{align}
\begin{pmatrix}
X_{-2}'\\ X_{-1}'\\ X_{\hat{+}}'\\ X_1'\\ X_2'\\ X_{\hat{-}}'
\end{pmatrix}
&=e^{-ib^{\hat{-}}\mathfrak{K}^{\hat{-}}}
\begin{pmatrix}
X_{-2}\\ X_{-1}\\ X_{\hat{+}}\\ X_1\\ X_2\\ X_{\hat{-}}
\end{pmatrix}
=\begin{pmatrix}
1 & 0 & 0 & 0 & 0 & 0\\
- \mathbb{C} (b^{\hat{-}})^{2} & 1 & 0 & 0 & 0 & \sqrt{2} b^{\hat{-}}\\
\sqrt{2} \mathbb{S} b^{\hat{-}} & 0 & 1 & 0 & 0 & 0\\
0 & 0 & 0 & 1 & 0 & 0\\
0 & 0 & 0 & 0 & 1 & 0\\
- \sqrt{2} \mathbb{C} b^{\hat{-}} & 0 & 0 & 0 & 0 & 1
\end{pmatrix}
\begin{pmatrix}
X_{-2}\\ X_{-1}\\ X_{\hat{+}}\\ X_1\\ X_2\\ X_{\hat{-}}
\end{pmatrix}\\
&=\begin{pmatrix}
X_{-2}\\
X_{-1}+\sqrt{2}b^{\hat{-}}X_{\hat{-}}
-\mathbb{C}(b^{\hat{-}})^2X_{-2}\\
X_{\hat{+}}+\sqrt{2}\mathbb{S}b^{\hat{-}}X_{-2}\\
X_1\\
X_2\\
X_{\hat{-}}-\sqrt{2}\mathbb{C}b^{\hat{-}}X_{-2}
\end{pmatrix}.
\end{align}
\begin{align}
x_{\hat{+}}'
&=-\frac{1}{\sqrt{2}}\frac{X_{\hat{+}}'}{X_{-1}'}
=-\frac{1}{\sqrt{2}}
\frac{X_{\hat{+}}/X_{-1}+\sqrt{2}\mathbb{S}b^{\hat{-}}X_{-2}/X_{-1}}
{1+\sqrt{2}b^{\hat{-}}X_{\hat{-}}/X_{-1}
 -\mathbb{C}(b^{\hat{-}})^2X_{-2}/X_{-1}}\nonumber\\
&=\frac{x_{\hat{+}}-\mathbb{S}b^{\hat{-}}x^2}
{1-2b^{\hat{-}}x_{\hat{-}}-\mathbb{C}(b^{\hat{-}})^2x^2}, \nonumber\\
x_{\hat{-}}'
&=-\frac{1}{\sqrt{2}}\frac{X_{\hat{-}}'}{X_{-1}'}
=-\frac{1}{\sqrt{2}}
\frac{X_{\hat{-}}/X_{-1}-\sqrt{2}\mathbb{C}b^{\hat{-}}X_{-2}/X_{-1}}
{1+\sqrt{2}b^{\hat{-}}X_{\hat{-}}/X_{-1}
 -\mathbb{C}(b^{\hat{-}})^2X_{-2}/X_{-1}}\nonumber\\
&=\frac{x_{\hat{-}}+\mathbb{C}b^{\hat{-}}x^2}
{1-2b^{\hat{-}}x_{\hat{-}}-\mathbb{C}(b^{\hat{-}})^2x^2}.
\end{align}
The $b^{\hat{-}}x^2$ terms cancel when the index is raised:
\begin{align}
x^{\hat{+}\prime}
&=\frac{x^{\hat{+}}}
{1+2\mathbb{C}b^{\hat{-}}x^{\hat{-}}
 -2\mathbb{S}b^{\hat{-}}x^{\hat{+}}
 -\mathbb{C}(b^{\hat{-}})^2x^2}.
\end{align}

\subsection{Translations}

\subsubsection{\texorpdfstring{$P_{\hat{+}}$}{P hat plus}}

The transformation of the interpolating projective space-time vector by
$P_{\hat{+}}$ is
\begin{align}
\begin{pmatrix}
X_{-2}'\\ X_{-1}'\\ X_{\hat{+}}'\\ X_1'\\ X_2'\\ X_{\hat{-}}'
\end{pmatrix}
&=e^{-ia^{\hat{+}}P_{\hat{+}}}
\begin{pmatrix}
X_{-2}\\ X_{-1}\\ X_{\hat{+}}\\ X_1\\ X_2\\ X_{\hat{-}}
\end{pmatrix}
=\begin{pmatrix}
1 & \mathbb{C} (a^{\hat{+}})^{2} & - \sqrt{2} a^{\hat{+}} & 0 & 0 & 0\\
0 & 1 & 0 & 0 & 0 & 0\\
0 & - \sqrt{2} \mathbb{C} a^{\hat{+}} & 1 & 0 & 0 & 0\\
0 & 0 & 0 & 1 & 0 & 0\\
0 & 0 & 0 & 0 & 1 & 0\\
0 & - \sqrt{2} \mathbb{S} a^{\hat{+}} & 0 & 0 & 0 & 1
\end{pmatrix}
\begin{pmatrix}
X_{-2}\\ X_{-1}\\ X_{\hat{+}}\\ X_1\\ X_2\\ X_{\hat{-}}
\end{pmatrix}\\
&=\begin{pmatrix}
X_{-2}+\mathbb{C}(a^{\hat{+}})^2X_{-1}
-\sqrt{2}a^{\hat{+}}X_{\hat{+}}\\
X_{-1}\\
X_{\hat{+}}-\sqrt{2}\mathbb{C}a^{\hat{+}}X_{-1}\\
X_1\\ X_2\\
X_{\hat{-}}-\sqrt{2}\mathbb{S}a^{\hat{+}}X_{-1}
\end{pmatrix}.
\end{align}
Therefore,
\begin{align}
x_{\hat{+}}'
&=-\frac{1}{\sqrt{2}}\frac{X_{\hat{+}}-\sqrt{2}\mathbb{C}a^{\hat{+}}X_{-1}}{X_{-1}}
=x_{\hat{+}}+\mathbb{C}a^{\hat{+}}, \nonumber\\
x_{\hat{-}}'
&=-\frac{1}{\sqrt{2}}\frac{X_{\hat{-}}-\sqrt{2}\mathbb{S}a^{\hat{+}}X_{-1}}{X_{-1}}
=x_{\hat{-}}+\mathbb{S}a^{\hat{+}}, \nonumber\\
x^{\hat{+}\prime}
&=\mathbb{C}x_{\hat{+}}'+\mathbb{S}x_{\hat{-}}'
=x^{\hat{+}}+a^{\hat{+}}.
\end{align}

\subsubsection{\texorpdfstring{$P_{\hat{-}}$}{P hat minus}}

The transformation of the interpolating projective space-time vector by
$P_{\hat{-}}$ is
\begin{align}
\begin{pmatrix}
X_{-2}'\\ X_{-1}'\\ X_{\hat{+}}'\\ X_1'\\ X_2'\\ X_{\hat{-}}'
\end{pmatrix}
&=e^{-ia^{\hat{-}}P_{\hat{-}}}
\begin{pmatrix}
X_{-2}\\ X_{-1}\\ X_{\hat{+}}\\ X_1\\ X_2\\ X_{\hat{-}}
\end{pmatrix}
=\begin{pmatrix}
1 & - \mathbb{C} (a^{\hat{-}})^{2} & 0 & 0 & 0 & - \sqrt{2} a^{\hat{-}}\\
0 & 1 & 0 & 0 & 0 & 0\\
0 & - \sqrt{2} \mathbb{S} a^{\hat{-}} & 1 & 0 & 0 & 0\\
0 & 0 & 0 & 1 & 0 & 0\\
0 & 0 & 0 & 0 & 1 & 0\\
0 & \sqrt{2} \mathbb{C} a^{\hat{-}} & 0 & 0 & 0 & 1
\end{pmatrix}
\begin{pmatrix}
X_{-2}\\ X_{-1}\\ X_{\hat{+}}\\ X_1\\ X_2\\ X_{\hat{-}}
\end{pmatrix}\\
&=\begin{pmatrix}
X_{-2}-\mathbb{C}(a^{\hat{-}})^2X_{-1}
-\sqrt{2}a^{\hat{-}}X_{\hat{-}}\\
X_{-1}\\
X_{\hat{+}}-\sqrt{2}\mathbb{S}a^{\hat{-}}X_{-1}\\
X_1\\ X_2\\
X_{\hat{-}}+\sqrt{2}\mathbb{C}a^{\hat{-}}X_{-1}
\end{pmatrix}.
\end{align}
so that
\begin{align}
x_{\hat{+}}'
&=-\frac{1}{\sqrt{2}}\frac{X_{\hat{+}}-\sqrt{2}\mathbb{S}a^{\hat{-}}X_{-1}}{X_{-1}}
=x_{\hat{+}}+\mathbb{S}a^{\hat{-}}, \nonumber\\
x_{\hat{-}}'
&=-\frac{1}{\sqrt{2}}\frac{X_{\hat{-}}+\sqrt{2}\mathbb{C}a^{\hat{-}}X_{-1}}{X_{-1}}
=x_{\hat{-}}-\mathbb{C}a^{\hat{-}}, \nonumber\\
x^{\hat{+}\prime}
&=\mathbb{C}x_{\hat{+}}'+\mathbb{S}x_{\hat{-}}'
=x^{\hat{+}}.
\end{align}

\subsection{Interpolating dilatation transformations}

\subsubsection{\texorpdfstring{$D_{\hat{+}}$}{D hat plus}}

Let $u=\alpha_{\hat{+}}\sin\delta$. The transformation of the interpolating projective
space-time vector by $D_{\hat{+}}$ is
\begin{align}
\begin{pmatrix}
X_{-2}'\\ X_{-1}'\\ X_{\hat{+}}'\\ X_1'\\ X_2'\\ X_{\hat{-}}'
\end{pmatrix}
&=e^{-i\alpha_{\hat{+}}D_{\hat{+}}}
\begin{pmatrix}
X_{-2}\\ X_{-1}\\ X_{\hat{+}}\\ X_1\\ X_2\\ X_{\hat{-}}
\end{pmatrix}\\
&=\begin{pmatrix}
e^{-\alpha_{\hat{+}}\cos\delta} & 0 & 0 & 0 & 0 & 0\\
0 & e^{\alpha_{\hat{+}}\cos\delta} & 0 & 0 & 0 & 0\\
0 & 0 & \cosh u-\mathbb{S}\sinh u & 0 & 0 & \mathbb{C}\sinh u\\
0 & 0 & 0 & 1 & 0 & 0\\
0 & 0 & 0 & 0 & 1 & 0\\
0 & 0 & \mathbb{C}\sinh u & 0 & 0 & \cosh u+\mathbb{S}\sinh u
\end{pmatrix}
\begin{pmatrix}
X_{-2}\\ X_{-1}\\ X_{\hat{+}}\\ X_1\\ X_2\\ X_{\hat{-}}
\end{pmatrix}\\
&=\begin{pmatrix}
e^{-\alpha_{\hat{+}}\cos\delta}X_{-2}\\
e^{\alpha_{\hat{+}}\cos\delta}X_{-1}\\
\left(\cosh u-\mathbb{S}\sinh u\right)X_{\hat{+}}
 +\mathbb{C}\sinh u\,X_{\hat{-}}\\
X_1\\ X_2\\
\mathbb{C}\sinh u\,X_{\hat{+}}
 +\left(\cosh u+\mathbb{S}\sinh u\right)X_{\hat{-}}
\end{pmatrix}.
\end{align}
It follows that
\begin{align}
x_{\hat{+}}'&=e^{-\alpha_{\hat{+}}\cos\delta}
\left[\left(\cosh u-\mathbb{S}\sinh u\right)x_{\hat{+}}
 +\mathbb{C}\sinh u\,x_{\hat{-}}\right], \nonumber\\
x_{\hat{-}}'&=e^{-\alpha_{\hat{+}}\cos\delta}
\left[\mathbb{C}\sinh u\,x_{\hat{+}}
 +\left(\cosh u+\mathbb{S}\sinh u\right)x_{\hat{-}}\right].
\end{align}
Raising the index yields
\begin{align}
x^{\hat{+}\prime}
=e^{-\alpha_{\hat{+}}\cos\delta}
\left[\left(\cosh u+\mathbb{S}\sinh u\right)x^{\hat{+}}
 -\mathbb{C}\sinh u\,x^{\hat{-}}\right].
\end{align}

\subsubsection{\texorpdfstring{$D_{\hat{-}}$}{D hat minus}}

Let $v=\alpha_{\hat{-}}\cos\delta$. The transformation of the interpolating projective
space-time vector by $D_{\hat{-}}$ is
\begin{align}
\begin{pmatrix}
X_{-2}'\\ X_{-1}'\\ X_{\hat{+}}'\\ X_1'\\ X_2'\\ X_{\hat{-}}'
\end{pmatrix}
&=e^{-i\alpha_{\hat{-}}D_{\hat{-}}}
\begin{pmatrix}
X_{-2}\\ X_{-1}\\ X_{\hat{+}}\\ X_1\\ X_2\\ X_{\hat{-}}
\end{pmatrix}\\
&=\begin{pmatrix}
e^{-\alpha_{\hat{-}}\sin\delta} & 0 & 0 & 0 & 0 & 0\\
0 & e^{\alpha_{\hat{-}}\sin\delta} & 0 & 0 & 0 & 0\\
0 & 0 & \cosh v+\mathbb{S}\sinh v & 0 & 0 & -\mathbb{C}\sinh v\\
0 & 0 & 0 & 1 & 0 & 0\\
0 & 0 & 0 & 0 & 1 & 0\\
0 & 0 & -\mathbb{C}\sinh v & 0 & 0 & \cosh v-\mathbb{S}\sinh v
\end{pmatrix}
\begin{pmatrix}
X_{-2}\\ X_{-1}\\ X_{\hat{+}}\\ X_1\\ X_2\\ X_{\hat{-}}
\end{pmatrix}\\
&=\begin{pmatrix}
e^{-\alpha_{\hat{-}}\sin\delta}X_{-2}\\
e^{\alpha_{\hat{-}}\sin\delta}X_{-1}\\
\left(\cosh v+\mathbb{S}\sinh v\right)X_{\hat{+}}
 -\mathbb{C}\sinh v\,X_{\hat{-}}\\
X_1\\ X_2\\
-\mathbb{C}\sinh v\,X_{\hat{+}}
 +\left(\cosh v-\mathbb{S}\sinh v\right)X_{\hat{-}}
\end{pmatrix}.
\end{align}
Thus,
\begin{align}
x_{\hat{+}}'&=e^{-\alpha_{\hat{-}}\sin\delta}
\left[\left(\cosh v+\mathbb{S}\sinh v\right)x_{\hat{+}}
 -\mathbb{C}\sinh v\,x_{\hat{-}}\right], \nonumber\\
x_{\hat{-}}'&=e^{-\alpha_{\hat{-}}\sin\delta}
\left[-\mathbb{C}\sinh v\,x_{\hat{+}}
 +\left(\cosh v-\mathbb{S}\sinh v\right)x_{\hat{-}}\right],
\end{align}
and hence
\begin{align}
x^{\hat{+}\prime}
=e^{-\alpha_{\hat{-}}\sin\delta}
\left[\left(\cosh v-\mathbb{S}\sinh v\right)x^{\hat{+}}
 +\mathbb{C}\sinh v\,x^{\hat{-}}\right].
\end{align}

\subsection{Rotation and transverse special conformal transformations}

\subsubsection{\texorpdfstring{$J^{3}$}{J3}}

The transformation of the interpolating projective space-time vector by
$J^3$ is
\begin{align}
\begin{pmatrix}
X_{-2}'\\ X_{-1}'\\ X_{\hat{+}}'\\ X_1'\\ X_2'\\ X_{\hat{-}}'
\end{pmatrix}
&=e^{-i\theta J^3}
\begin{pmatrix}
X_{-2}\\ X_{-1}\\ X_{\hat{+}}\\ X_1\\ X_2\\ X_{\hat{-}}
\end{pmatrix}
=\begin{pmatrix}
1 & 0 & 0 & 0 & 0 & 0\\
0 & 1 & 0 & 0 & 0 & 0\\
0 & 0 & 1 & 0 & 0 & 0\\
0 & 0 & 0 & \cos\theta & -\sin\theta & 0\\
0 & 0 & 0 & \sin\theta & \cos\theta & 0\\
0 & 0 & 0 & 0 & 0 & 1
\end{pmatrix}
\begin{pmatrix}
X_{-2}\\ X_{-1}\\ X_{\hat{+}}\\ X_1\\ X_2\\ X_{\hat{-}}
\end{pmatrix}\\
&=\begin{pmatrix}
X_{-2}\\
X_{-1}\\
X_{\hat{+}}\\
\cos\theta X_1-\sin\theta X_2\\
\sin\theta X_1+\cos\theta X_2\\
X_{\hat{-}}
\end{pmatrix}.
\end{align}
and
\begin{align}
x_{\hat{+}}'=-\frac{1}{\sqrt{2}}\frac{X_{\hat{+}}}{X_{-1}}=x_{\hat{+}},\qquad
x_{\hat{-}}'=-\frac{1}{\sqrt{2}}\frac{X_{\hat{-}}}{X_{-1}}=x_{\hat{-}},\qquad
x^{\hat{+}\prime}=x^{\hat{+}}.
\end{align}

\subsubsection{\texorpdfstring{$\mathfrak{K}_{1}$}{K1}}

The transformation of the interpolating projective space-time vector by
$\mathfrak{K}_1$ is
\begin{align}
\begin{pmatrix}
X_{-2}'\\ X_{-1}'\\ X_{\hat{+}}'\\ X_1'\\ X_2'\\ X_{\hat{-}}'
\end{pmatrix}
&=e^{-ib^1\mathfrak{K}_1}
\begin{pmatrix}
X_{-2}\\ X_{-1}\\ X_{\hat{+}}\\ X_1\\ X_2\\ X_{\hat{-}}
\end{pmatrix}
=\begin{pmatrix}
1 & 0 & 0 & 0 & 0 & 0\\
- (b^{1})^{2} & 1 & 0 & \sqrt{2} b^{1} & 0 & 0\\
0 & 0 & 1 & 0 & 0 & 0\\
- \sqrt{2} b^{1} & 0 & 0 & 1 & 0 & 0\\
0 & 0 & 0 & 0 & 1 & 0\\
0 & 0 & 0 & 0 & 0 & 1
\end{pmatrix}
\begin{pmatrix}
X_{-2}\\ X_{-1}\\ X_{\hat{+}}\\ X_1\\ X_2\\ X_{\hat{-}}
\end{pmatrix}\\
&=\begin{pmatrix}
X_{-2}\\
X_{-1}+\sqrt{2}b^1X_1-(b^1)^2X_{-2}\\
X_{\hat{+}}\\
X_1-\sqrt{2}b^1X_{-2}\\
X_2\\
X_{\hat{-}}
\end{pmatrix}.
\end{align}
\begin{align}
x_{\hat{+}}'
&=-\frac{1}{\sqrt{2}}\frac{X_{\hat{+}}'}{X_{-1}'}
=-\frac{1}{\sqrt{2}}
\frac{X_{\hat{+}}/X_{-1}}
{1+\sqrt{2}b^1X_1/X_{-1}-(b^1)^2X_{-2}/X_{-1}}\nonumber\\
&=\frac{x_{\hat{+}}}{1-2b^1x_1-(b^1)^2x^2}, \nonumber\\
x_{\hat{-}}'
&=-\frac{1}{\sqrt{2}}\frac{X_{\hat{-}}'}{X_{-1}'}
=-\frac{1}{\sqrt{2}}
\frac{X_{\hat{-}}/X_{-1}}
{1+\sqrt{2}b^1X_1/X_{-1}-(b^1)^2X_{-2}/X_{-1}}\nonumber\\
&=\frac{x_{\hat{-}}}{1-2b^1x_1-(b^1)^2x^2}, \nonumber\\
x^{\hat{+}\prime}
&=\mathbb{C}x_{\hat{+}}'+\mathbb{S}x_{\hat{-}}'
=\frac{x^{\hat{+}}}{1-2b^1x_1-(b^1)^2x^2}.
\end{align}

\subsubsection{\texorpdfstring{$\mathfrak{K}_{2}$}{K2}}

The transformation of the interpolating projective space-time vector by
$\mathfrak{K}_2$ is
\begin{align}
\begin{pmatrix}
X_{-2}'\\ X_{-1}'\\ X_{\hat{+}}'\\ X_1'\\ X_2'\\ X_{\hat{-}}'
\end{pmatrix}
&=e^{-ib^2\mathfrak{K}_2}
\begin{pmatrix}
X_{-2}\\ X_{-1}\\ X_{\hat{+}}\\ X_1\\ X_2\\ X_{\hat{-}}
\end{pmatrix}
=\begin{pmatrix}
1 & 0 & 0 & 0 & 0 & 0\\
- (b^{2})^{2} & 1 & 0 & 0 & \sqrt{2} b^{2} & 0\\
0 & 0 & 1 & 0 & 0 & 0\\
0 & 0 & 0 & 1 & 0 & 0\\
- \sqrt{2} b^{2} & 0 & 0 & 0 & 1 & 0\\
0 & 0 & 0 & 0 & 0 & 1
\end{pmatrix}
\begin{pmatrix}
X_{-2}\\ X_{-1}\\ X_{\hat{+}}\\ X_1\\ X_2\\ X_{\hat{-}}
\end{pmatrix}\\
&=\begin{pmatrix}
X_{-2}\\
X_{-1}+\sqrt{2}b^2X_2-(b^2)^2X_{-2}\\
X_{\hat{+}}\\
X_1\\
X_2-\sqrt{2}b^2X_{-2}\\
X_{\hat{-}}
\end{pmatrix}.
\end{align}
\begin{align}
x_{\hat{+}}'
&=-\frac{1}{\sqrt{2}}\frac{X_{\hat{+}}'}{X_{-1}'}
=-\frac{1}{\sqrt{2}}
\frac{X_{\hat{+}}/X_{-1}}
{1+\sqrt{2}b^2X_2/X_{-1}-(b^2)^2X_{-2}/X_{-1}}\nonumber\\
&=\frac{x_{\hat{+}}}{1-2b^2x_2-(b^2)^2x^2}, \nonumber\\
x_{\hat{-}}'
&=-\frac{1}{\sqrt{2}}\frac{X_{\hat{-}}'}{X_{-1}'}
=-\frac{1}{\sqrt{2}}
\frac{X_{\hat{-}}/X_{-1}}
{1+\sqrt{2}b^2X_2/X_{-1}-(b^2)^2X_{-2}/X_{-1}}\nonumber\\
&=\frac{x_{\hat{-}}}{1-2b^2x_2-(b^2)^2x^2}, \nonumber\\
x^{\hat{+}\prime}
&=\mathbb{C}x_{\hat{+}}'+\mathbb{S}x_{\hat{-}}'
=\frac{x^{\hat{+}}}{1-2b^2x_2-(b^2)^2x^2}.
\end{align}

\subsection{Transverse translations}

\subsubsection{\texorpdfstring{$P_{1}$}{P1}}

The transformation of the interpolating projective space-time vector by
$P_1$ is
\begin{align}
\begin{pmatrix}
X_{-2}'\\ X_{-1}'\\ X_{\hat{+}}'\\ X_1'\\ X_2'\\ X_{\hat{-}}'
\end{pmatrix}
&=e^{-ia^1P_1}
\begin{pmatrix}
X_{-2}\\ X_{-1}\\ X_{\hat{+}}\\ X_1\\ X_2\\ X_{\hat{-}}
\end{pmatrix}
=\begin{pmatrix}
1 & - (a^{1})^{2} & 0 & - \sqrt{2} a^{1} & 0 & 0\\
0 & 1 & 0 & 0 & 0 & 0\\
0 & 0 & 1 & 0 & 0 & 0\\
0 & \sqrt{2} a^{1} & 0 & 1 & 0 & 0\\
0 & 0 & 0 & 0 & 1 & 0\\
0 & 0 & 0 & 0 & 0 & 1
\end{pmatrix}
\begin{pmatrix}
X_{-2}\\ X_{-1}\\ X_{\hat{+}}\\ X_1\\ X_2\\ X_{\hat{-}}
\end{pmatrix}\\
&=\begin{pmatrix}
X_{-2}-(a^1)^2X_{-1}-\sqrt{2}a^1X_1\\
X_{-1}\\ X_{\hat{+}}\\ X_1+\sqrt{2}a^1X_{-1}\\ X_2\\ X_{\hat{-}}
\end{pmatrix}.
\end{align}
\begin{align}
x_{\hat{+}}'=-\frac{1}{\sqrt{2}}\frac{X_{\hat{+}}}{X_{-1}}=x_{\hat{+}},\qquad
x_{\hat{-}}'=-\frac{1}{\sqrt{2}}\frac{X_{\hat{-}}}{X_{-1}}=x_{\hat{-}},\qquad
x^{\hat{+}\prime}=x^{\hat{+}}.
\end{align}

\subsubsection{\texorpdfstring{$P_{2}$}{P2}}

The transformation of the interpolating projective space-time vector by
$P_2$ is
\begin{align}
\begin{pmatrix}
X_{-2}'\\ X_{-1}'\\ X_{\hat{+}}'\\ X_1'\\ X_2'\\ X_{\hat{-}}'
\end{pmatrix}
&=e^{-ia^2P_2}
\begin{pmatrix}
X_{-2}\\ X_{-1}\\ X_{\hat{+}}\\ X_1\\ X_2\\ X_{\hat{-}}
\end{pmatrix}
=\begin{pmatrix}
1 & - (a^{2})^{2} & 0 & 0 & - \sqrt{2} a^{2} & 0\\
0 & 1 & 0 & 0 & 0 & 0\\
0 & 0 & 1 & 0 & 0 & 0\\
0 & 0 & 0 & 1 & 0 & 0\\
0 & \sqrt{2} a^{2} & 0 & 0 & 1 & 0\\
0 & 0 & 0 & 0 & 0 & 1
\end{pmatrix}
\begin{pmatrix}
X_{-2}\\ X_{-1}\\ X_{\hat{+}}\\ X_1\\ X_2\\ X_{\hat{-}}
\end{pmatrix}\\
&=\begin{pmatrix}
X_{-2}-(a^2)^2X_{-1}-\sqrt{2}a^2X_2\\
X_{-1}\\ X_{\hat{+}}\\ X_1\\ X_2+\sqrt{2}a^2X_{-1}\\ X_{\hat{-}}
\end{pmatrix}.
\end{align}
\begin{align}
x_{\hat{+}}'=-\frac{1}{\sqrt{2}}\frac{X_{\hat{+}}}{X_{-1}}=x_{\hat{+}},\qquad
x_{\hat{-}}'=-\frac{1}{\sqrt{2}}\frac{X_{\hat{-}}}{X_{-1}}=x_{\hat{-}},\qquad
x^{\hat{+}\prime}=x^{\hat{+}}.
\end{align}

\subsection{Transverse interpolating boosts}

\subsubsection{\texorpdfstring{$\mathcal{D}^{1}$}{D1}}

Set $q_1=\sqrt{\mathbb{C}}\eta^{1}$. The transformation of the interpolating
projective space-time vector by $\mathcal{D}^{1}$ is
\begin{align}
\begin{pmatrix}
X_{-2}'\\ X_{-1}'\\ X_{\hat{+}}'\\ X_1'\\ X_2'\\ X_{\hat{-}}'
\end{pmatrix}
&=e^{-i\eta^{1}\mathcal{D}^{1}}
\begin{pmatrix}
X_{-2}\\ X_{-1}\\ X_{\hat{+}}\\ X_1\\ X_2\\ X_{\hat{-}}
\end{pmatrix}
=\begin{pmatrix}
1 & 0 & 0 & 0 & 0 & 0\\
0 & 1 & 0 & 0 & 0 & 0\\
0 & 0 & \cosh q_1 & \sqrt{\mathbb{C}}\sinh q_1 & 0 & 0\\
0 & 0 & \dfrac{\sinh q_1}{\sqrt{\mathbb{C}}} & \cosh q_1 & 0 & 0\\
0 & 0 & 0 & 0 & 1 & 0\\
0 & 0 & \dfrac{\mathbb{S}(\cosh q_1-1)}{\mathbb{C}} & \dfrac{\mathbb{S}\sinh q_1}{\sqrt{\mathbb{C}}} & 0 & 1
\end{pmatrix}
\begin{pmatrix}
X_{-2}\\ X_{-1}\\ X_{\hat{+}}\\ X_1\\ X_2\\ X_{\hat{-}}
\end{pmatrix}\\
&=\begin{pmatrix}
X_{-2}\\ X_{-1}\\
\cosh q_1\,X_{\hat{+}}+\sqrt{\mathbb{C}}\sinh q_1\,X_1\\
\dfrac{\sinh q_1}{\sqrt{\mathbb{C}}}X_{\hat{+}}+\cosh q_1\,X_1\\
X_2\\
X_{\hat{-}}+\dfrac{\mathbb{S}(\cosh q_1-1)}{\mathbb{C}}X_{\hat{+}}
+\dfrac{\mathbb{S}\sinh q_1}{\sqrt{\mathbb{C}}}X_1
\end{pmatrix}.
\end{align}
Therefore,
\begin{align}
x_{\hat{+}}'&=\cosh q_1\,x_{\hat{+}}
 +\sqrt{\mathbb{C}}\sinh q_1\,x_1, \nonumber\\
x_{\hat{-}}'&=x_{\hat{-}}
 +\frac{\mathbb{S}(\cosh q_1-1)}{\mathbb{C}}x_{\hat{+}}
 +\frac{\mathbb{S}\sinh q_1}{\sqrt{\mathbb{C}}}x_1.
\end{align}
Using the interpolating metric to raise the index gives
\begin{align}
x^{\hat{+}\prime}
&=\cosh q_1\,x^{\hat{+}}
 +\frac{\mathbb{S}}{\mathbb{C}}\left(\cosh q_1-1\right)x^{\hat{-}}
 +\frac{\sinh q_1}{\sqrt{\mathbb{C}}}x_1.
\end{align}

\subsubsection{\texorpdfstring{$\mathcal{D}^{2}$}{D2}}

With $q_2=\sqrt{\mathbb{C}}\eta^{2}$, the transformation of the interpolating
projective space-time vector by $\mathcal{D}^{2}$ is
\begin{align}
\begin{pmatrix}
X_{-2}'\\ X_{-1}'\\ X_{\hat{+}}'\\ X_1'\\ X_2'\\ X_{\hat{-}}'
\end{pmatrix}
&=e^{-i\eta^{2}\mathcal{D}^{2}}
\begin{pmatrix}
X_{-2}\\ X_{-1}\\ X_{\hat{+}}\\ X_1\\ X_2\\ X_{\hat{-}}
\end{pmatrix}
=\begin{pmatrix}
1 & 0 & 0 & 0 & 0 & 0\\
0 & 1 & 0 & 0 & 0 & 0\\
0 & 0 & \cosh q_2 & 0 & \sqrt{\mathbb{C}}\sinh q_2 & 0\\
0 & 0 & 0 & 1 & 0 & 0\\
0 & 0 & \dfrac{\sinh q_2}{\sqrt{\mathbb{C}}} & 0 & \cosh q_2 & 0\\
0 & 0 & \dfrac{\mathbb{S}(\cosh q_2-1)}{\mathbb{C}} & 0 & \dfrac{\mathbb{S}\sinh q_2}{\sqrt{\mathbb{C}}} & 1
\end{pmatrix}
\begin{pmatrix}
X_{-2}\\ X_{-1}\\ X_{\hat{+}}\\ X_1\\ X_2\\ X_{\hat{-}}
\end{pmatrix}\\
&=\begin{pmatrix}
X_{-2}\\ X_{-1}\\
\cosh q_2\,X_{\hat{+}}+\sqrt{\mathbb{C}}\sinh q_2\,X_2\\
X_1\\
\dfrac{\sinh q_2}{\sqrt{\mathbb{C}}}X_{\hat{+}}+\cosh q_2\,X_2\\
X_{\hat{-}}+\dfrac{\mathbb{S}(\cosh q_2-1)}{\mathbb{C}}X_{\hat{+}}
+\dfrac{\mathbb{S}\sinh q_2}{\sqrt{\mathbb{C}}}X_2
\end{pmatrix}.
\end{align}
Hence,
\begin{align}
x_{\hat{+}}'&=\cosh q_2\,x_{\hat{+}}
 +\sqrt{\mathbb{C}}\sinh q_2\,x_2, \nonumber\\
x_{\hat{-}}'&=x_{\hat{-}}
 +\frac{\mathbb{S}(\cosh q_2-1)}{\mathbb{C}}x_{\hat{+}}
 +\frac{\mathbb{S}\sinh q_2}{\sqrt{\mathbb{C}}}x_2,
\end{align}
and
\begin{align}
x^{\hat{+}\prime}
&=\cosh q_2\,x^{\hat{+}}
 +\frac{\mathbb{S}}{\mathbb{C}}\left(\cosh q_2-1\right)x^{\hat{-}}
 +\frac{\sinh q_2}{\sqrt{\mathbb{C}}}x_2.
\end{align}

\subsection{Transverse interpolating rotations}

\subsubsection{\texorpdfstring{$\mathcal{K}^{1}$}{K1}}

Let $r_1=\sqrt{\mathbb{C}}\phi^{1}$. The transformation of the interpolating
projective space-time vector by $\mathcal{K}^{1}$ is
\begin{align}
\begin{pmatrix}
X_{-2}'\\ X_{-1}'\\ X_{\hat{+}}'\\ X_1'\\ X_2'\\ X_{\hat{-}}'
\end{pmatrix}
&=e^{-i\phi^{1}\mathcal{K}^{1}}
\begin{pmatrix}
X_{-2}\\ X_{-1}\\ X_{\hat{+}}\\ X_1\\ X_2\\ X_{\hat{-}}
\end{pmatrix}
=\begin{pmatrix}
1 & 0 & 0 & 0 & 0 & 0\\
0 & 1 & 0 & 0 & 0 & 0\\
0 & 0 & 1 & \dfrac{\mathbb{S}\sin r_1}{\sqrt{\mathbb{C}}} & 0 & \dfrac{\mathbb{S}(1-\cos r_1)}{\mathbb{C}}\\
0 & 0 & 0 & \cos r_1 & 0 & \dfrac{\sin r_1}{\sqrt{\mathbb{C}}}\\
0 & 0 & 0 & 0 & 1 & 0\\
0 & 0 & 0 & -\sqrt{\mathbb{C}}\sin r_1 & 0 & \cos r_1
\end{pmatrix}
\begin{pmatrix}
X_{-2}\\ X_{-1}\\ X_{\hat{+}}\\ X_1\\ X_2\\ X_{\hat{-}}
\end{pmatrix}\\
&=\begin{pmatrix}
X_{-2}\\ X_{-1}\\
X_{\hat{+}}+\dfrac{\mathbb{S}\sin r_1}{\sqrt{\mathbb{C}}}X_1
+\dfrac{\mathbb{S}(1-\cos r_1)}{\mathbb{C}}X_{\hat{-}}\\
\cos r_1\,X_1+\dfrac{\sin r_1}{\sqrt{\mathbb{C}}}X_{\hat{-}}\\
X_2\\
-\sqrt{\mathbb{C}}\sin r_1\,X_1+\cos r_1\,X_{\hat{-}}
\end{pmatrix}.
\end{align}
The transformed lower coordinates are
\begin{align}
x_{\hat{+}}'&=x_{\hat{+}}
 +\frac{\mathbb{S}\sin r_1}{\sqrt{\mathbb{C}}}x_1
 +\frac{\mathbb{S}(1-\cos r_1)}{\mathbb{C}}x_{\hat{-}}, \nonumber\\
x_{\hat{-}}'&=-\sqrt{\mathbb{C}}\sin r_1\,x_1
 +\cos r_1\,x_{\hat{-}}.
\end{align}
The transverse and $x_{\hat{-}}$ contributions cancel after raising the
index:
\begin{align}
x^{\hat{+}\prime}
&=\mathbb{C}x_{\hat{+}}'+\mathbb{S}x_{\hat{-}}'
=\mathbb{C}x_{\hat{+}}+\mathbb{S}x_{\hat{-}}
=x^{\hat{+}}.
\end{align}

\subsubsection{\texorpdfstring{$\mathcal{K}^{2}$}{K2}}

For $r_2=\sqrt{\mathbb{C}}\phi^{2}$, the transformation of the interpolating
projective space-time vector by $\mathcal{K}^{2}$ is
\begin{align}
\begin{pmatrix}
X_{-2}'\\ X_{-1}'\\ X_{\hat{+}}'\\ X_1'\\ X_2'\\ X_{\hat{-}}'
\end{pmatrix}
&=e^{-i\phi^{2}\mathcal{K}^{2}}
\begin{pmatrix}
X_{-2}\\ X_{-1}\\ X_{\hat{+}}\\ X_1\\ X_2\\ X_{\hat{-}}
\end{pmatrix}
=\begin{pmatrix}
1 & 0 & 0 & 0 & 0 & 0\\
0 & 1 & 0 & 0 & 0 & 0\\
0 & 0 & 1 & 0 & \dfrac{\mathbb{S}\sin r_2}{\sqrt{\mathbb{C}}} & \dfrac{\mathbb{S}(1-\cos r_2)}{\mathbb{C}}\\
0 & 0 & 0 & 1 & 0 & 0\\
0 & 0 & 0 & 0 & \cos r_2 & \dfrac{\sin r_2}{\sqrt{\mathbb{C}}}\\
0 & 0 & 0 & 0 & -\sqrt{\mathbb{C}}\sin r_2 & \cos r_2
\end{pmatrix}
\begin{pmatrix}
X_{-2}\\ X_{-1}\\ X_{\hat{+}}\\ X_1\\ X_2\\ X_{\hat{-}}
\end{pmatrix}\\
&=\begin{pmatrix}
X_{-2}\\ X_{-1}\\
X_{\hat{+}}+\dfrac{\mathbb{S}\sin r_2}{\sqrt{\mathbb{C}}}X_2
+\dfrac{\mathbb{S}(1-\cos r_2)}{\mathbb{C}}X_{\hat{-}}\\
X_1\\
\cos r_2\,X_2+\dfrac{\sin r_2}{\sqrt{\mathbb{C}}}X_{\hat{-}}\\
-\sqrt{\mathbb{C}}\sin r_2\,X_2+\cos r_2\,X_{\hat{-}}
\end{pmatrix}.
\end{align}
\begin{align}
x_{\hat{+}}'&=x_{\hat{+}}
 +\frac{\mathbb{S}\sin r_2}{\sqrt{\mathbb{C}}}x_2
 +\frac{\mathbb{S}(1-\cos r_2)}{\mathbb{C}}x_{\hat{-}}, \nonumber\\
x_{\hat{-}}'&=-\sqrt{\mathbb{C}}\sin r_2\,x_2
 +\cos r_2\,x_{\hat{-}}, \nonumber\\
x^{\hat{+}\prime}
&=\mathbb{C}x_{\hat{+}}'+\mathbb{S}x_{\hat{-}}'
=x^{\hat{+}}.
\end{align}

\reftitle{References}
\bibliography{BibTeXList}

\end{document}
